% paralleles.tex — f162v, vv. 1–42 : loci paralleli (registre établi à partir des collectes vérifiées, voir edition/CONVENTIONS.md § 4)
\lp{1 ὦ Ζεῦ … φανήσεται]}{Ar.~\textit{Nub.}~2–4 ὦ Ζεῦ βασιλεῦ … οὐδέποθ᾽ ἡμέρα γενήσεται; / καὶ μὴν πάλαι γ᾽ ἀλεκτρυόνος ἤκουσ᾽ ἐγώ ; Ar.~\textit{Pax}~58 καί φησιν, ὦ Ζεῦ τί ποτε βουλεύει ποιεῖν ; Ar.~\textit{Pax}~62 ὦ Ζεῦ τί δρασείεις ποθ᾽ ἡμῶν τὸν λεών ; Ar.~\textit{Lys.}~476 ὦ Ζεῦ τί ποτε χρησόμεθα τοῖσδε τοῖς κνωδάλοις ; Ar.~\textit{Thesm.}~71–72 ὦ Ζεῦ τί δρᾶσαι διανοεῖ με τήμερον; / νὴ τοὺς θεοὺς ἐγὼ πυθέσθαι βούλομαι ; Ar.~\textit{Plut.}~1–2 ὡς ἀργαλέον πρᾶγμ᾽ ἐστὶν ὦ Ζεῦ καὶ θεοὶ / δοῦλον γενέσθαι παραφρονοῦντος δεσπότου}
\lp{1 ἆρά ποτε]}{Ar.~\textit{Pax}~1087 ἆρα φενακίζων ποτ᾽ Ἀθηναίους ἔτι παύσει ; [Eur.] \textit{Rhes.}~360–362 ἆρά ποτ᾽ αὖθις ἁ παλαιὰ Τροΐα / τοὺς προπότας παναμερεύ- / σει θιάσους ἐρώτων}
\lp{1 χελιδὼν … φανήσεται]}{Ar.~fr.~617~K.-A. πυθοῦ χελιδὼν πηνίκ᾽ ἄττα φαίνεται ; Chionid. fr.~8~Kock πυθοῦ χελιδὼν πηνίκ᾽ ἄττα φαίνεται ; Σ \textit{Thesm.}~1 τῷ δὲ ἔαρι χελιδόνες φαίνονται … οἷον πότε ἀπαλλαγήσομαι τοῦ κακοῦ τούτου ; Eust.~148.46 πύθου, χελιδὼν πηνίκ᾽ ἄττα φαίνεται ; Hsch.~π~4311 πύθου χελιδόνος· παροιμιῶδες, διότι ὀδυρτικὸν τὸ ζῷον· θρηνεῖ γὰρ ἡ χελιδών ; Phot.~π~1519 πυθοῦ χελιδόνος· παροιμία ἀπό τινος Χελιδόνος θεολόγου ἀνδρὸς … οἱ δὲ ὅτι θρηνητικὸν τὸ ζῷον· οἱ δὲ ὅτι τὸ ἔαρ προσημαίνει ; Suid.~π~3136 Πυθοῦ Χελιδόνος: παροιμία ἀπό τινος Χελιδόνος, θεολόγου καὶ τερατοσκόπου ; EM~167.46 ὁ κωμικὸς δὲ χρονικὸν αὐτὸ παραλαμβάνει· πύθου χελιδὼν τηνίκ᾽ ἄττα γίνεται ; Harp. s.v. ἄττα ἐνιαχοῦ δὲ παρέλκει τὸ ἄττα. [πυθοῦ χελιδὼν πηνίκ᾽ ἄττα φαίνεται.] ; Apostol.~15.14 Πύθου Χελιδόνος: ἀντὶ τοῦ, ἐγκαίρως ποίει τὰ σαυτοῦ ; App. Prov.~4.63 Πύθου χελιδόνος: διότι ὀδυρτικὸν τὸ ζῶον}
\lp{1 χελιδών]}{Eur.~fr.~773~N.² (Phaéthon) μέλπει δ᾽ ἐν δένδρεσι λεπτὰν / ἀηδὼν ἁρμονίαν / ὀρθρευομένα γόοις / Ἴτυν Ἴτυν πολύθρηνον ; Hes.~\textit{Op.}~568–569 τὸν δὲ μέτ᾽ ὀρθογόη Πανδιονὶς ὦρτο χελιδὼν / ἐς φάος ἀνθρώποις, ἔαρος νέον ἱσταμένοιο ; Sapph. fr.~135 V. (ap. Heph.~12.2) τί με Πανδιονὶς ὤραννα χελιδών ; Anacreont.~9.8–10~Bergk τί μευ καλῶν ὀνείρων / ὑπορθρίαισι φωναῖς / ἀφήρπασας Βάθυλλον ; AP~5.237.1–3 (Agathias) εὖτε δ᾽ ἐπέλθῃ / ὄρθρος … ἀμφιπεριτρύζουσι χελιδόνες ; AP~6.160.1–2 (Antipater de Sidon) κερκίδα τὰν ὀρθρινά, χελιδονίδων ἅμα φωνᾷ, / μελπομέναν, ἱστῶν Παλλάδος ἀλκυόνα ; AP~6.247.1 (Philippe) κερκίδας ὀρθρολάλοισι χελιδόσιν εἰκελοφώνους ; Hsch.~α~1226 ἀδωνηίς· ἡ χελιδών. καὶ ἡ θριδακίνη ; Hsch.~ο~1189 ὀρθογόη· ἡ χελιδών ; Hsch.~ω~302 ὦ ᾽ράνα χελίδων· ὀροφή ; Phot.~ο~469 ὀρθογόη· χελιδών ; Suid.~ο~587 Ὀρθρογόη: χελιδών. λέγεται καὶ ὀρθρόλαλος}
\lp{1 χελιδών (printemps)]}{Ar.~\textit{Eq.}~419 σκέψασθε παῖδες· οὐχ ὁρᾶθ᾽; ὥρα νέα, χελιδών ; Ar.~\textit{Pax}~800–801 ὅταν ἠρινὰ μὲν φωνῇ χελιδὼν / ἑζομένη κελαδῇ ; Ar.~\textit{Av.}~714–715 εἶτα χελιδών, / ὅτε χρὴ χλαῖναν πωλεῖν ἤδη καὶ ληδάριόν τι πρίασθαι ; Ar.~\textit{Av.}~1416–1417 ἐς θοἰμάτιον τὸ σκόλιον ᾄδειν μοι δοκεῖ, / δεῖσθαι δ᾽ ἔοικεν οὐκ ὀλίγων χελιδόνων ; Philem. fr.~208~Kock ἡ μὲν χελιδὼν τὸ θέρος, ὦ γύναι, λαλεῖ ; Simon. fr.~74~Bergk ἄγγελε κλυτὰ ἔαρος ἁδυόδμου, / κυανέα χελιδοῖ ; \textit{Carm.\,pop.} ap.~Ath.~8.360c ἦλθ᾽, ἦλθε χελιδὼν / καλὰς ὥρας ἄγουσα, / καλοὺς ἐνιαυτούς ; \textit{Pervig.\,Ven.}~89–90 Quando ver venit meum? / Quando fiam uti chelidon ut tacere desinam? ; Arist.~\textit{EN}~1098a18 μία γὰρ χελιδὼν ἔαρ οὐ ποιεῖ, οὐδὲ μία ἡμέρα ; Σ Ar.~\textit{Av.}~1417 μετείληφε δὲ τὸν νοῦν τῆς παροιμίας, μία χελιδὼν ἔαρ οὐ ποιεῖ ; Σ Ar.~\textit{Eq.}~419 νέα χελιδών: Παροιμιῶδες … δοκεῖ γάρ πως ἅμα τῷ ἔαρι φαίνεσθαι ἡ χελιδών ; Hsch.~μ~1318 μία χελιδών· παροιμιῶδες τοῦτο, ὅτι μία χελιδὼν ἔαρ οὐ ποιεῖ ; Hsch.~χ~324 χελιδονισταί· οἱ τῇ χελιδόνι ἀγείροντες ; Zen.~5.12 Μία χελιδὼν ἔαρ οὐ ποιεῖ ; Tosi~705 Μία χελιδὼν ἔαρ οὐ ποιεῖ}
\lp{1 χελιδών (babil)]}{Ar.~\textit{Ran.}~93 χελιδόνων μουσεῖα, λωβηταὶ τέχνης ; Ar.~\textit{Ran.}~678–681 Κλεοφῶντος, ἐφ᾽ οὗ δὴ χείλεσιν ἀμφιλάλοις / δεινὸν ἐπιβρέμεται / Θρῃκία χελιδών ; Aesch.~\textit{Ag.}~1050 ἀλλ᾽ εἴπερ ἐστι μὴ χελιδόνος δίκην ; Thphr.~\textit{Char.}~7.9 οὐκ ἂν σιωπήσειεν, οὐδ᾽ εἰ τῶν χελιδόνων δόξειεν εἶναι λαλίστερος ; Porph.~\textit{VP}~42 μηδὲ χελιδόνας ἐν οἰκίᾳ δέχεσθαι, τουτέστι λάλους ἀνθρώπους ; Σ Ar.~\textit{Ran.}~93 παροιμία ἐπὶ τῶν πολυλόγων καὶ ἐπαχθῶν ταττομένη ; Hsch.~τ~517 τερετίζοντα· λαλοῦντα. ἐκ μεταφορᾶς τῆς χελιδόνος ; Hsch.~τ~993 τιτυβίζει· ὡς χελιδὼν φωνεῖ ; Hsch.~χ~325 χελιδόνος δίκην· τοὺς βαρβάρους χελιδόσιν ἀπεικάζουσι διὰ τὴν ἀσύνετον λαλιάν ; Hsch.~χ~327 χελιδόνων μουσεῖα· ὡς βάρβαρα καὶ ἀσύνετα ποιούντων τῶν τραγικῶν ; Tosi~2058 Μηδὲ χελιδόνας ἐν οἰκίᾳ δέχεσθαι}
\lp{2 ἀπολεῖ μ᾽ ἀλοῶν ἅνθρωπος ἐξ ἑωθινοῦ]}{\textit{Test.} Suid.~α~1330 Ἀλοῶν, ἔξωθεν ἐν κύκλῳ περιάγων, ὡς οἱ ἐν ταῖς ἅλωσιν. ἢ τύπτων … Ἀριστοφάνης· ἀπολεῖ μ᾽ ἀλοῶν ἄνθρωπος ἐξ ἑωθινοῦ}
\lp{2–4 ἀπολεῖ μ᾽ ἀλοῶν … ποῖ μ᾽ ἄγεις]}{Ar.~\textit{Av.}~3–6 τί ὦ πόνηρ᾽ ἄνω κάτω πλανύττομεν; … ἀπολούμεθ᾽ ἄλλως τὴν ὁδὸν προφορουμένω}
\lp{2 ἀπολεῖ μ᾽]}{Ar.~\textit{Ach.}~470 ἀπολεῖς μ᾽. ἰδού σοι. φροῦδά μοι τὰ δράματα ; Ar.~\textit{Thesm.}~1073 ἀπολεῖς μ᾽, ὦ γραῦ, στωμυλλομένη}
\lp{2 ἀλοῶν]}{\textit{Test.} Synag. Σᵇ α~984~Cunningham ἀλοῶν· ἀντὶ τοῦ περιάγων, ὡς οἱ ἀλοῶντες βόες. οὕτως Ἀριστοφάνης (Th. 2) ; Ar.~\textit{Ran.}~149–150 ἢ μητέρ᾽ ἠλόασεν, ἢ πατρὸς γνάθον / ἐπάταξεν ; Ar.~fr.~932~K.-A. (661~Kock) ἀλοᾶν χρὴ τὰς γνάθους ; Eub. fr.~15.5~Kock κατηλόηται γαστρὸς ἐν μέσῳ κύκλος ; Alex. fr.~147.2~Kock ἄνω κάτω τε περιπατοῦσ᾽ ὥσπερ Πλάτων ; Soph.~fr.~19~N.² κέστρᾳ σιδηρᾷ πλευρὰ καὶ κατὰ ῥάχιν / ἠλοῆσαι πλεῖον ; Theoc.~22.128 μέχρι συνηλοίησε παρήια ; Herod.~2.34 κἠμὲ τὸν ξεῖνον / οὐ[δεὶ]ς πολίτης ἠλόησεν ; Babr.~98.15–16 τὸν δ᾽ ἕκαστος ἠλοία / ῥοπάλῳ τις ἢ λίθῳ τις ἐκ χερὸς παίων ; Pl.~\textit{Phdr.}~230d–e σὺ ἐμοὶ λόγους οὕτω προτείνων ἐν βιβλίοις τήν τε Ἀττικὴν φαίνῃ περιάξειν ἅπασαν ; Pl.~\textit{Thg.}~124a τῶν φυτευόντων καὶ σπειρόντων καὶ ἁλοώντων ; X.~\textit{Cyr.}~7.1.31 τοὺς δὲ πίπτοντας κατηλόων καὶ αὐτοὺς καὶ ὅπλα καὶ ἵπποις καὶ τροχοῖς ; Aeschin.~2.140 ὅτε αὐτῶν τοὺς ὁμήρους λαβόντες Φωκεῖς κατηλόησαν ; Luc.~\textit{Icar.}~15 σκύφῳ χρυσῷ τὴν ὀφρὺν κατηλοημένος ; Σ \textit{Thesm.}~2 Ἔξωθεν ἐν κύκλῳ περιάγων ὡς οἱ ἐν ταῖς ἅλωσι. καὶ τὸ τύπτειν ἀλοᾶν λέγουσιν ; Σ Luc.~\textit{Icar.}~15 ἀλοᾶν γὰρ τὸ συντρίβειν ἢ τύπτειν … ἀφ᾽ οὗ καὶ πατραλοίας ; Hsch.~α~3221 ἀλοία· ἔτυπτεν (Ι 568), ἀφ᾽ οὗ καὶ ὁ πατροτύπτης πατραλοίας ; Hsch.~α~3233 ἀλοᾷ· τύπτει ; Hsch.~α~3266 ἀλοῶν· πλανῶν. καὶ τύπτων ; Suid.~α~1305 τὸ δὲ συγκόψαι πληγαῖς, ἀλοῆσαι ἐν τῷ η ; Synag. Σᵇ α~986~Cunningham ἀλοᾶν· τύπτειν, βάλλειν. καὶ τὸ θρύπτειν ἐπὶ τῆς ἅλω τὸν σῖτον}
\lp{2 ἐξ ἑωθινοῦ]}{Ar.~\textit{Ach.}~19–20 ὡς νῦν, ὁπότ᾽ οὔσης κυρίας ἐκκλησίας / ἑωθινῆς ἔρημος ἡ πνὺξ αὑτηί ; Ar.~\textit{Thesm.}~376 ἐκκλησίαν ποιεῖν ἕωθεν τῇ μέσῃ ; Ar.~\textit{Eccl.}~84–85 ἡκκλησία δ᾽, εἰς ἣν παρεσκευάσμεθα / ἡμεῖς βαδίζειν, ἐξ ἕω γενήσεται ; Pherecr. fr.~90~Kock ζητῶ περιέρρων αὐτὸν ἐξ ἑωθινοῦ ; Alex. fr.~257.4–5~Kock εὐθὺς ἐξ ἑωθινοῦ / ἕστηκεν ἐλθών ; Pl.~\textit{Lg.}~722c ἐξ ἑωθινοῦ μεσημβρία τε γέγονε ; Pl.~\textit{Phdr.}~227a συχνὸν γὰρ ἐκεῖ διέτριψα χρόνον καθήμενος ἐξ ἑωθινοῦ ; Pl.~\textit{Phdr.}~228b καὶ τοῦτο δρῶν ἐξ ἑωθινοῦ καθήμενος ἀπειπὼν εἰς περίπατον ᾔει ; Pl.~\textit{Smp.}~220c Σωκράτης ἐξ ἑωθινοῦ φροντίζων τι ἕστηκε ; Hsch.~ε~7732 ἑωθινός· ὀρθρινός}
\lp{3–4 οἷόν τε … πυθέσθαι]}{Ar.~\textit{Plut.}~18–19 ἐγὼ μὲν οὖν οὐκ ἔσθ᾽ ὅπως σιγήσομαι, / ἢν μὴ φράσῃς ὅ τι τῷδ᾽ ἀκολουθοῦμέν ποτε}
\lp{3 οἷόν τε]}{Ar.~\textit{Eq.}~74 ἀλλ᾽ οὐχ οἷόν τε τὸν Παφλαγόν᾽ οὐδὲν λαθεῖν ; Ar.~\textit{Nub.}~217 ἀλλ᾽ οὐχ οἷόν τε ; Σ \textit{Thesm.}~3 οἷόν τε : δυνατὸν, ἐρωτηματικῶς}
\lp{3 τὸν σπλῆνα … ἐκβαλεῖν]}{Nicopho fr.~12~Kock (ap.~Ath.~3.80b) πυρετὸς εὐθέως ἥκει τρέχων οὐκ ἄξιος τριωβόλου κᾆθ᾽ οὗτος ἐπιπεσὼν ἐμεῖν ποιεῖ χολήν ; Plaut.~\textit{Cas.}~414–415 perii, cor lienosum, opinor, habeo, iam dudum salit, / de labore pectus tundit ; Plaut.~\textit{Merc.}~124 perii, seditionem facit lien, occupat praecordia ; Herod.~3.70 δότω τις εἰς τὴν χεῖρα πρὶν χολὴ⟨ν⟩ βῆξαι ; Plu.~\textit{Dem.}~6.2 καχεξίαν τινὰ σπληνὸς ἀμυνόμενον δρόμοις μακροῖς χρῆσθαι ; Plu.~Mor.~831c (De vit. aere al.~8) ἐμοῦντος τοῦ ἑτέρου καὶ λέγοντος τὰ σπλάγχν᾽ ἐκβάλλειν ; Plin.~\textit{NH}~11.205 peculiare cursus inpedimentum aliquando in eo, quam ob rem minuitur cursorum laborantibus ; Plin.~\textit{NH}~26.83 Equisaetum … lienes cursorum extinguit decocta fictili novo ad tertias ; Hsch.~δ~2733 δώραξ· σπλὴν ὑπὸ Μακεδόνων ; Hsch.~σ~1531 σπλῆνες· οἱ ὑπὸ τῷ λέβητι ἕλικες. καὶ παρὰ τοῖς ἰατροῖς τὰ ἐπιμήκη ὀθόνια ; Greg. Cypr.~2.35 Ἔνεστι καὶ μύρμηκι χολή: ὅτι μηδὲ τοῦ μικροῦ δεῖ καταφρονεῖν ; Apostol.~8.25 Ἔχει καὶ οἱ μυῖα σπλῆνα: Ἔχει καὶ χολὴν ὁ μύρμηξ: ἐπὶ τῶν εὐκαταφρονήτων}
\lp{3 κομιδῇ]}{Ar.~\textit{Pax}~820 ἔγωγέ τοι πεπόνηκα κομιδῇ τὼ σκέλει ; Ar.~\textit{Plut.}~1086 ἀλλ᾽ ἔστι κομιδῇ τρὺξ παλαιὰ καὶ σαπρά ; Hsch.~κ~3450 κομιδή· ἄγαν, πάνυ, σφόδρα ; Suid.~κ~1999 Κομιδῇ: λίαν, παντελῶς, ὁλοσχερῶς, εἰλικρινῶς, παντάπασι, τελείως}
\lp{4 παρὰ σοῦ πυθέσθαι]}{Ar.~\textit{Thesm.}~72 νὴ τοὺς θεοὺς ἐγὼ πυθέσθαι βούλομαι ; Gorg.~\textit{Pal.}~26 Βουλοίμην δ᾽ ἂν παρὰ σοῦ πυθέσθαι}
\lp{4 ποῖ μ᾽ ἄγεις]}{Alex. fr.~99~Kock ποῖ δή μ᾽ ἄγεις διὰ τῶν κύκλων ; Luc.~\textit{DMar.}~8.3 Ἄνθρωπε, ποῖ με συναρπάσας ἄγεις}
\lp{4 ωὖριπίδη]}{Ar.~\textit{Ach.}~404 Εὐριπίδη, Εὐριπίδιον ; Ar.~\textit{Nub.}~221–222 ὦ Σώκρατες, / ὦ Σωκρατίδιον}
\lp{5–6 ἀλλ᾽ οὐκ ἀκούειν … ὄψει παρεστώς]}{Plaut.~\textit{Truc.}~489 pluris est oculatus testis unus quam auriti decem ; Eur.~\textit{Med.}~652–653 εἴδομεν, οὐκ ἐξ ἑτέρων / μῦθον ἔχω φράσασθαι ; Eur.~\textit{Or.}~81 Ἑλένη, τί σοι λέγοιμ᾽ ἂν ἅ γε παροῦσ᾽ ὁρᾷς ; Aesch.~\textit{Pers.}~266–267 καὶ μὴν παρών γε κοὐ λόγους ἄλλων κλύων, / Πέρσαι, φράσαιμ᾽ ἂν οἷ᾽ ἐπορσύνθη κακά ; Soph.~\textit{OT}~6–7 ἁγὼ δικαιῶν μὴ παρ᾽ ἀγγέλων, τέκνα, / ἄλλων ἀκούειν αὐτὸς ὧδ᾽ ἐλήλυθα ; Soph.~\textit{Trach.}~746–747 αὐτὸς βαρεῖαν ξυμφορὰν ἐν ὄμμασιν / πατρὸς δεδορκὼς κοὐ κατὰ γλῶσσαν κλύων ; Emp.~31 B~3.9–10~DK (ap. S.E. M.~7.125) ἀλλ᾽ ἄγ᾽ ἄθρει πάσῃ παλάμῃ, πῇ δῆλον ἕκαστον, / μήτε τιν᾽ ὄψιν ἔχων πίστει πλέον ἢ κατ᾽ ἀκουήν ἢ ἀκοὴν ἐρίδουπον ὑπὲρ τρανώματα γλώσσης ; Hor.~\textit{Ars}~180–181 segnius inritant animos demissa per aurem / quam quae sunt oculis subiecta fidelibus ; Heraclit.~22 B~101a~DK (Plb.~12.27.1) ὀφθαλμοὶ γὰρ τῶν ὤτων ἀκριβέστεροι μάρτυρες ; Hdt.~1.8.2 ὦτα γὰρ τυγχάνει ἀνθρώποισι ἐόντα ἀπιστότερα ὀφθαλμῶν ; Antipho Soph.~87 B~44~DK (P. Oxy.~1364, A col.~2.31–3.5) νενομοθέτηται γὰρ ἐπί τε τοῖς ὀφθαλμοῖς, ἃ δεῖ αὐτοὺς ὁρᾶν καὶ ἃ οὐ δεῖ· καὶ ἐπὶ τοῖς ὠσίν, ἃ δεῖ αὐτὰ ἀκούειν καὶ ἃ οὐ δεῖ ; Th.~1.73.2 ὧν ἀκοαὶ μᾶλλον λόγων μάρτυρες ἢ ὄψις τῶν ἀκουσομένων ; Pl.~\textit{Tht.}~201b περὶ ὧν ἰδόντι μόνον ἔστιν εἰδέναι, ἄλλως δὲ μή, ταῦτα τότε ἐξ ἀκοῆς κρίνοντες ; Luc.~\textit{Hist.\,Conscr.}~29 Ὦτα ὀφθαλμῶν ἀπιστότερα. γράφω τοίνυν ἃ εἶδον, οὐχ ἃ ἤκουσα ; Luc.~\textit{Salt.}~78 Ἡροδότῳ μὲν οὖν τὰ δι᾽ ὀμμάτων φαινόμενα πιστότερα εἶναι τῶν ὤτων δοκεῖ ; Sen.~\textit{Ep.}~6.5 homines amplius oculis quam auribus credunt ; Σ \textit{Thesm.}~5 ὁ μὲν τραγικώτερον καὶ ὑψηλότερον φράζει, ὁ δὲ ταπεινότερον ἢ δεῖ ἀκούει ; Apostol.~18.71 Ὠτίων πιστότεροι ὀφθαλμοί: ἐπὶ τῶν ἀλόγως τὰ μείω τοῖς κρείττοσι παραβαλλομένων ; Tosi~1943 Ὀφθαλμοὶ… τῶν ὤτων ἀκριβέστεροι μάρτυρες}
\lp{5 αὐτίκα]}{Ar.~\textit{Eccl.}~1061 αὐτοῦ τι δρῶντα πυρρὸν ὄψει μ᾽ αὐτίκα ; Hsch.~α~8392 αὐτίκα· Ἀριστοφάνης μὲν ἐπὶ τοῦ ἐνεστῶτος· Ὅμηρος δὲ ἐπὶ μόνου τοῦ νῦν}
\lp{6 πῶς λέγεις; αὖθις φράσον]}{Ar.~\textit{Av.}~323 πῶς λέγεις ; Ar.~\textit{Ran.}~515 πῶς λέγεις; ὀρχηστρίδες ; Ar.~\textit{Plut.}~268 ὦ χρυσὸν ἀγγείλας ἐπῶν πῶς φῄς; πάλιν φράσον μοι ; Eur.~\textit{Hel.}~471 πῶς φῄς; τίν᾽ εἶπας μῦθον; αὖθίς μοι φράσον ; Aesch.~\textit{Cho.}~767 ἦ πῶς; λέγ᾽ αὖθις, ὡς μάθω σαφέστερον ; Soph.~\textit{OT}~361 οὐχ ὥστε γ᾽ εἰπεῖν γνωστόν· ἀλλ᾽ αὖθις φράσον ; Soph.~\textit{Phil.}~341–342 τοιγαροῦν τὸ σὸν φράσον / αὖθις πάλιν μοι πρᾶγμ᾽, ὅτῳ σ᾽ ἐνύβρισαν ; Soph.~\textit{Phil.}~1407 πῶς λέγεις ; Hsch.~α~8269 αὖθις· πάλιν, ἐξ ἀρχῆς, εὐθύς}
\lp{7 οὐχ ἅ γ᾽ ἂν μέλλῃς ὁρᾶν]}{Ar.~\textit{Ran.}~960 ἐξ ὧν γ᾽ ἂν ἐξηλεγχόμην ; Eur.~\textit{Ba.}~924 νῦν δ᾽ ὁρᾷς ἃ χρή σ᾽ ὁρᾶν ; Eur.~\textit{Hipp.}~345a πῶς ἂν σύ μοι λέξειας ἁμὲ χρὴ λέγειν ; Eur.~\textit{Ion}~558 νῦν ὁρᾷς ἃ χρή σ᾽ ὁρᾶν ; Eur.~\textit{Or.}~81 Ἑλένη, τί σοι λέγοιμ᾽ ἂν ἅ γε παροῦσ᾽ ὁρᾷς ; Eur.~\textit{Phoen.}~713 μῶν νεάζων οὐχ ὁρᾷς ἃ χρή σ᾽ ὁρᾶν}
\lp{7–8 οὐχ ἅ γ᾽ ἂν (ἄγαν)]}{Eur.~\textit{Hipp.}~264–265 οὕτω τὸ λίαν ἧσσον ἐπαινῶ / τοῦ μηδὲν ἄγαν ; Soph.~\textit{OC}~1695 μηδὲν ἄγαν φλέγεσθον ; Pl.~\textit{Chrm.}~165a ὃ δή μοι δοκοῦσιν παθεῖν καὶ οἱ τὰ ὕστερον γράμματα ἀναθέντες, τό τε μηδὲν ἄγαν ; Pl.~\textit{Prt.}~343b γνῶθι σαυτόν καὶ μηδὲν ἄγαν ; Arist.~\textit{Rh.}~2.21.13 τὸ γνῶθι σαυτὸν καὶ τὸ μηδὲν ἄγαν ; Hsch.~α~303 ἄγαν· πάνυ, πολύ, μεγάλως, λίαν ; Greg. Cypr. L~2.79 Μηδὲν ἄγαν: ἑνὸς τῶν ἑπτὰ σοφῶν τὸ ἀπόφθεγμα, ὅμοιον τῷ· μηδὲν ὑπὲρ τὸ μέτρον ; Macar.~5.90 Μηδὲν ἄγαν: ἑνὸς τῶν ἑπτὰ σοφῶν τὸ ἀπόφθεγμα ; Tosi~1526 Ne quid nimis}
\lp{7–10 ὁρᾶν … ἀκούειν]}{Eur.~fr.~413.3~N.² (Ino) ὁρᾶν θ᾽ ἃ δεῖ με κοὐχ ὁρᾶν ἃ μὴ πρέπει}
\lp{8 οὐδ᾽ ἆρ᾽ ὁρᾶν δεῖ μ᾽]}{Ar.~\textit{Nub.}~1252 οὐκ ἄρ᾽ ἀποδώσεις ; Ar.~\textit{Av.}~91 οὐκ ἆρ᾽ ἀφῆκας ; Ar.~\textit{Eccl.}~462 οὐδὲ στένειν τὸν ὄρθρον ἔτι πρᾶγμ᾽ ἆρά μοι ; Ar.~\textit{Eccl.}~672 οὐδὲ κυβεύσουσ᾽ ἆρ᾽ ἄνθρωποι ; Soph.~\textit{Aj.}~1238 οὐκ ἆρ᾽ Ἀχαιοῖς ἄνδρες εἰσὶ πλὴν ὅδε ; Soph.~\textit{Phil.}~106 οὐκ ἆρ᾽ ἐκείνῳ γ᾽ οὐδὲ προσμῖξαι θρασύ ; Soph.~\textit{Phil.}~114 οὐκ ἆρ᾽ ὁ πέρσων, ὡς ἐφάσκετ᾽, εἴμ᾽ ἐγώ}
\lp{9 πῶς μοι παραινεῖς; δεξιῶς μέντοι λέγεις]}{Ar.~\textit{Ran.}~1169 εὖ νὴ τὸν Ἑρμῆν· ὅ τι λέγεις δ᾽ οὐ μανθάνω ; Ar.~\textit{Ran.}~1444–1445 πῶς; οὐ μανθάνω. / ἀμαθέστερόν πως εἰπὲ καὶ σαφέστερον}
\lp{9 παραινεῖς]}{Ar.~\textit{Nub.}~1483 ὀρθῶς παραινεῖς οὐκ ἐῶν δικορραφεῖν ; Aesch.~\textit{Cho.}~903 κρίνω σὲ νικᾶν, καὶ παραινεῖς μοι καλῶς ; Soph.~\textit{El.}~1472 ἀλλ᾽ εὖ παραινεῖς κἀπιπείσομαι ; Hsch.~π~519 παραίνεσις· συμβουλή, νουθεσία, σωφρονισμός}
\lp{9 δεξιῶς]}{Ar.~\textit{Eq.}~1377 σοφός γ᾽ ὁ Φαίαξ δεξιῶς τ᾽ οὐκ ἀπέθανεν ; Ar.~\textit{Nub.}~418–419 ὅπερ εἰκὸς δεξιὸν ἄνδρα, / νικᾶν πράττων καὶ βουλεύων καὶ τῇ γλώττῃ πολεμίζων ; Ar.~\textit{Nub.}~521 ὡς ὑμᾶς ἡγούμενος εἶναι θεατὰς δεξιοὺς ; Ar.~\textit{Ran.}~1370 ἐπίπονοί γ᾽ οἱ δεξιοί ; Ar.~\textit{Eccl.}~241 εὖ γ᾽ ὦ γλυκυτάτη Πραξαγόρα καὶ δεξιῶς ; Hsch.~δ~637 δεξιός· συνετός, ἀγαθός, καλός. ἐπιδέξιος. εὔθετος}
\lp{9 μέντοι λέγεις]}{Ar.~\textit{Lys.}~1095 νὴ τὸν Δί᾽ εὖ μέντοι λέγεις ; Ar.~\textit{Plut.}~1052b εὖ μέντοι λέγει ; Pl.~\textit{Lg.}~666a ἀναγκαιότατα μέντοι λέγεις}
\lp{10 οὔτ᾽ ἀκούειν οὔθ᾽ ὁρᾶν]}{Epich. fr.~12 (23 B~12~DK) νοῦς ὁρῆι καὶ νοῦς ἀκούει· τἆλλα κωφὰ καὶ τυφλά ; Men. Mon.~39~Meineke Ἃ μὴ προσήκει μήτ᾽ ἄκουε μήθ᾽ ὅρα ; Aesch.~\textit{PV}~447–448 οἳ πρῶτα μὲν βλέποντες ἔβλεπον μάτην, / κλύοντες οὐκ ἤκουον ; Soph.~\textit{OT}~371 τυφλὸς τά τ᾽ ὦτα τόν τε νοῦν τά τ᾽ ὄμματ᾽ εἶ ; Critias fr. Sisyph. (S.E. M.~9.54) ὡς ἔστι δαίμων ἀφθίτῳ θάλλων βίῳ, / νόῳ τ᾽ ἀκούων καὶ βλέπων, φρονῶν τε ; Emp.~31 B~2~DK (D.L.~9.73) οὕτως οὔτ᾽ ἐπιδερκτὰ τάδ᾽ ἀνδράσιν οὔτ᾽ ἐπακουστὰ οὔτε νόῳ περιληπτά ; Heraclit.~22 B~34~DK ἀξύνετοι ἀκούσαντες κωφοῖσιν ἐοίκασι· φάτις αὐτοῖσιν μαρτυρεῖ παρεόντας ἀπεῖναι ; Pl.~\textit{Phd.}~65b οἱ ποιηταὶ ἡμῖν ἀεὶ θρυλοῦσιν, ὅτι οὔτ᾽ ἀκούομεν ἀκριβὲς οὐδὲν οὔτε ὁρῶμεν ; D.~25.89 τὸ τῆς παροιμίας, ὁρῶντας μὴ ὁρᾶν καὶ ἀκούοντας μὴ ἀκούειν ; \textit{Ev.\,Matt.}~13.13 βλέποντες οὐ βλέπουσιν καὶ ἀκούοντες οὐκ ἀκούουσιν οὐδὲ συνίουσιν ; Plu.~Mor.~13e (De lib. educ.~18) ὡς ἔνια τῶν πραττομένων ὁρῶντας μὴ ὁρᾶν καὶ μὴ ἀκούειν ἀκούοντας ; Luc.~\textit{Vit.\,Auct.}~27 ἡ ἀμαθία καὶ τὸ μήτε ἀκούειν μήτε ὁρᾶν ; Apostol.~12.13 Νοῦς ὁρᾷ, καὶ νοῦς ἀκούει}
\lp{11–12 χωρὶς … ἡ φύσις … τοῦ μήτ᾽ ἀκούειν μήθ᾽ ὁρᾶν]}{Eur.~\textit{Alc.}~528 χωρὶς τό τ᾽ εἶναι καὶ τὸ μὴ νομίζεται ; Soph.~\textit{OC}~808 χωρὶς τό τ᾽ εἰπεῖν πολλὰ καὶ τὰ καίρια ; Hp.~\textit{Nat.\,Hom.}~5 ἔπειτα κατὰ φύσιν τὰς ἰδέας κεχωρίσθαι ; Pl.~\textit{Phlb.}~44a εἴπερ χωρὶς τοῦ μὴ λυπεῖσθαι καὶ τοῦ χαίρειν ἡ φύσις ἑκατέρου. καὶ μὴν χωρίς γε ἦν ; Σ \textit{Thesm.}~11 ταῦτα τῶν φυσικῶν λόγων. χωρὶς, τοῦ ἀκούειν καὶ ὁρᾶν ; Σ \textit{Thesm.}~12 ἀντὶ τοῦ εἰπεῖν τοῦ ἀκούειν καὶ ὁρᾶν εἶπε τοῦ μήτ᾽ ἀκούειν μήθ᾽ ὁρᾶν ; Apostol.~18.47 Χωρὶς τό τ᾽ εἰπεῖν πολλὰ καὶ τὸ καίρια}
\lp{11 χωρίς]}{Ar.~\textit{Ran.}~1164 χωρὶς γὰρ ἄλλης συμφορᾶς ἐλήλυθεν ; Eur.~\textit{IA}~859 χωρὶς τἀμὰ καὶ Ἀγαμέμνονος ; Anaxag.~59 B~6~DK οὐδὲ χωρὶς ἔστιν εἶναι, ἀλλὰ πάντα παντὸς μοῖραν μετέχει ; Str.~12.4.4 παροιμιάζονται / χωρὶς τὰ Μυσῶν καὶ Φρυγῶν ὁρίσματα ; Plu.~Mor.~122c (De tuenda san.~1) χωρίς γὰρ ἔφη τὰ φιλοσόφων καὶ ἰατρῶν ὥσπερ τινῶν Μυσῶν καὶ Φρυγῶν ὁρίσματα ; Σ Luc.~31.21 Rabe ἀντὶ τοῦ ὅμοιον τῇ παροιμίᾳ τῇ διαγορευούσῃ χωρὶς τὰ Μυσῶν καὶ Φρυγῶν ὁρίσματα ; Hsch.~χ~876 χωρίς· δίχα. ἰδιαζόντως ; Greg. Cypr.~3.99 Χωρὶς τὰ Φρυγῶν καὶ Μυσῶν ὁρίσματα ; Apostol.~18.45 Χωρὶς τὰ Μυσῶν καὶ Φρυγῶν ὁρίσματα: ἐπὶ τῶν ἐναντίων}
\lp{12 εὖ ἴσθ᾽ ὅτι]}{Ar.~\textit{Pax}~374 εὖ ἴσθ᾽ ὅτι ; Ar.~\textit{Av.}~1408 ἀλλ᾽ οὖν ἔγωγ᾽ οὐ παύσομαι, τοῦτ᾽ ἴσθ᾽ ὅτι ; Ar.~\textit{Lys.}~154 σπονδὰς ποιήσαιντ᾽ ἂν ταχέως, εὖ οἶδ᾽ ὅτι ; Ar.~\textit{Plut.}~183 καὶ τῶν κακῶν καὶ τῶν ἀγαθῶν, εὖ ἴσθ᾽ ὅτι ; Ar.~\textit{Plut.}~889 μὰ τὸν Δί᾽ οὔκουν τῷ γε σῷ, σάφ᾽ ἴσθ᾽ ὅτι ; Hsch.~ε~6854 εὖ ἴσθι· καλῶς γίνωσκε}
\lp{13 πῶς χωρίς]}{Ar.~\textit{Eccl.}~126 πῶς καταγέλαστον}
\lp{13 οὕτω ταῦτα διεκρίθη τότε]}{Ar.~\textit{Vesp.}~1182 ἐκεῖνον ὡς οὕτω ποτ᾽ ἦν μῦς καὶ γαλῆ ; Pl.~\textit{Prt.}~320c ἦν γάρ ποτε χρόνος ὅτε θεοὶ μὲν ἦσαν, θνητὰ δὲ γένη οὐκ ἦν}
\lp{13 διεκρίθη]}{Epich. fr.~9 (23 B~9~DK) συνεκρίθη καὶ διεκρίθη κἀπῆλθεν ὅθεν ἦλθεν πάλιν ; Eur.~\textit{HF}~104 ἐξίσταται γὰρ πάντ᾽ ἀπ᾽ ἀλλήλων δίχα ; Eur.~fr.~484.2–3~N.² (Mélanippe la sage) ὡς οὐρανός τε γαῖά τ᾽ ἦν μορφὴ μία· / ἐπεὶ δ᾽ ἐχωρίσθησαν ἀλλήλων δίχα ; Eur.~fr.~839.13–14~N.² διακρινόμενον δ᾽ ἄλλο πρὸς ἄλλου / μορφὴν ἑτέραν ἀπέδειξεν ; Hes.~\textit{Th.}~535 καὶ γὰρ ὅτ᾽ ἐκρίνοντο θεοὶ θνητοί τ᾽ ἄνθρωποι ; A.R.~1.496–498 γαῖα καὶ οὐρανὸς ἠδὲ θάλασσα … νείκεος ἐξ ὀλοοῖο διέκριθεν ἀμφὶς ἕκαστα ; Ov.~\textit{Met.}~1.6 unus erat toto naturae vultus in orbe ; Anaxag.~59 B~13~DK καὶ ὅσον ἐκίνησεν ὁ νοῦς, πᾶν τοῦτο διεκρίθη ; D.S.~1.7.1 μίαν ἔχειν ἰδέαν οὐρανόν τε καὶ γῆν … μετὰ δὲ ταῦτα διαστάντων τῶν σωμάτων ἀπ᾽ ἀλλήλων ; D.S.~1.7.7 οὐδ᾽ Εὐριπίδης διαφωνεῖν τοῖς προειρημένοις, μαθητὴς ὢν Ἀναξαγόρου τοῦ φυσικοῦ ; Hsch.~δ~1082 διακρινθέντας· χωρισθέντας ; Hsch.~δ~1083 διακρινθεῖτε· διακριθείητε, διαλυθείητε ; Hsch.~δ~1085 διακρῖναι· διαχωρίσαι ; Hsch.~δ~1088 διακρίνουσιν· διαχωρίζουσιν. μεσιτεύουσι}
\lp{14 αἰθὴρ … διεχωρίζετο]}{Anaxag.~59 B~8~DK οὐ κεχώρισται ἀλλήλων τὰ ἐν τῷ ἑνὶ κόσμῳ οὐδὲ ἀποκέκοπται πελέκει ; Anaxag.~59 B~12~DK καὶ ἀποκρίνεται ἀπό τε τοῦ ἀραιοῦ τὸ πυκνὸν ; Emp.~31 A~49~DK ([Plut.] Plac.~2.6,~887c) Ἐμπεδοκλῆς τὸν μὲν αἰθέρα πρῶτον διακριθῆναι, δεύτερον δὲ τὸ πῦρ}
\lp{14 τὰ πρῶτα]}{Ar.~\textit{Av.}~693 Χάος ἦν καὶ Νὺξ Ἔρεβός τε μέλαν πρῶτον καὶ Τάρταρος εὐρύς ; \textit{Il.}~1.6 ἐξ οὗ δὴ τὰ πρῶτα διαστήτην ἐρίσαντε}
\lp{14 αἰθήρ]}{Ar.~\textit{Nub.}~265 λαμπρός τ᾽ Αἰθήρ ; Ar.~\textit{Nub.}~570 Αἰθέρα σεμνότατον βιοθρέμμονα πάντων ; Ar.~\textit{Thesm.}~272 ὄμνυμι τοίνυν αἰθέρ᾽ οἴκησιν Διός ; Ar.~\textit{Ran.}~100 αἰθέρα Διὸς δωμάτιον, ἢ χρόνου πόδα ; Ar.~\textit{Ran.}~892 αἰθὴρ ἐμὸν βόσκημα καὶ γλώσσης στρόφιγξ ; Eur.~\textit{Hel.}~34 εἴδωλον ἔμπνουν οὐρανοῦ ξυνθεῖσ᾽ ἄπο ; Eur.~\textit{Hel.}~583–584 καὶ τίς βλέποντα σώματ᾽ ἐξεργάζεται; / αἰθήρ, ὅθεν σὺ θεοπόνητ᾽ ἔχεις λέχη ; Eur.~\textit{Ion}~1079 ἀνεχόρευσεν αἰθήρ ; Eur.~\textit{Supp.}~533–534 πνεῦμα μὲν πρὸς αἰθέρα, / τὸ σῶμα δ᾽ ἐς γῆν ; Eur.~\textit{Tro.}~884 ὦ γῆς ὄχημα κἀπὶ γῆς ἔχων ἕδραν ; Eur.~fr.~330.1–2~N.² (Danaé) ἐς ταὐτὸν ἥκειν φημὶ ταῖς βροτῶν τύχαις / τόνδ᾽ ὃν καλοῦσιν αἰθέρ᾽, ᾧ τάδ᾽ ἔστι δή ; Eur.~fr.~487~N.² ὄμνυμι δ᾽ ἱερὸν αἰθέρ᾽, οἴκησιν Διός ; Eur.~fr.~839.1–2~N.² Γαῖα μεγίστη καὶ Διὸς Αἰθήρ, / ὃ μὲν ἀνθρώπων καὶ θεῶν γενέτωρ ; Eur.~fr.~877~N.² ἀλλ᾽ αἰθὴρ τίκτει σε, κόρα, / Ζεὺς ὃς ἀνθρώποις ὀνομάζεται ; Eur.~fr.~941~N.² ὁρᾷς τὸν ὑψοῦ τόνδ᾽ ἄπειρον αἰθέρα / καὶ γῆν πέριξ ἔχονθ᾽ ὑγραῖς ἐν ἀγκάλαις; / τοῦτον νόμιζε Ζῆνα, τόνδ᾽ ἡγοῦ θεόν ; Eur.~fr.~1023~N.² Αἰθέρα καὶ Γαῖαν πάντων γενέτειραν ἀείδω ; Hes.~\textit{Th.}~124 Νυκτὸς δ᾽ αὖτ᾽ Αἰθήρ τε καὶ Ἡμέρη ἐξεγένοντο ; Hsch.~α~1854 αἰθέριος ῥύμβος· οὐρανός ; Hsch.~α~1855 αἰθέρα· ἐμπυρισμόν, ἀπὸ τοῦ αἴθεσθαι. Εὐριπίδης Τρῳάσιν ; Hsch.~α~1860 αἰθήρ· ὁ ὑπὲρ τὰ νέφη τόπος ; Suid.~αι~121 Αἰθήρ: ὁ ἐν ὕψει ἀήρ, ὁ ἐπάνω τοῦ ἀέρος καιόμενος ἐκ τοῦ ἡλίου}
\lp{15 ζᾦ᾽ ἐν αὑτῷ ξυνετέκνου κινούμενα]}{Ar.~\textit{Pax}~832–833 οὐκ ἦν ἄρ᾽ οὐδ᾽ ἃ λέγουσι, κατὰ τὸν ἀέρα / ὡς ἀστέρες γιγνόμεθ᾽, ὅταν τις ἀποθάνῃ ; Eur.~\textit{Hel.}~1144–1146 σὺ Διὸς ἔφυς, ὦ Ἑλένα, θυγάτηρ· / πτανὸς γὰρ ἐν κόλποις σε Λή- / δας ἐτέκνωσε πατήρ ; Eur.~fr.~484.3–4~N.² ἐπεὶ δ᾽ ἐχωρίσθησαν ἀλλήλων δίχα, / τίκτουσι πάντα κἀνέδωκαν εἰς φάος ; Eur.~fr.~839.4–5~N.² παραδεξαμένη τίκτει θνητούς, / τίκτει βοτάνην φῦλά τε θηρῶν ; Eur.~fr.~898.7–10~N.² ἐρᾷ μὲν ὄμβρου γαῖ᾽, ὅταν ξηρὸν πέδον / ἄκαρπον αὐχμῷ νοτίδος ἐνδεῶς ἔχῃ· / ἐρᾷ δ᾽ ὁ σεμνὸς οὐρανὸς πληρούμενος / ὄμβρου πεσεῖν εἰς γαῖαν Ἀφροδίτης ὕπο· / ὅταν δὲ συμμιχθῆτον ἐς ταὐτὸν δύο, / φύουσιν ἡμῖν πάντα καὶ τρέφουσ ; Aesch.~fr. ap.~Ath.~13.600b (Danaïdes) ἐρᾷ μὲν ἁγνὸς οὐρανὸς τρῶσαι χθόνα ; Anaxag. ap. Thphr. HP~3.1.4 Ἀναξαγόρας μὲν τὸν ἀέρα πάντων φάσκων ἔχειν σπέρματα ; Pl.~\textit{Phd.}~96b ἆρ᾽ ἐπειδὰν τὸ θερμὸν καὶ τὸ ψυχρὸν σηπεδόνα τινὰ λάβῃ … τότε δὴ τὰ ζῷα συντρέφεται ; D.S.~1.7.3 γενέσθαι περὶ αὐτὰ σηπεδόνας ὑμέσι λεπτοῖς περιεχομένας ; Archel. ap. D.L.~2.17 γεννᾶσθαι δέ φησι τὰ ζῷα ἐκ θερμῆς τῆς γῆς ; Σ \textit{Thesm.}~15 κινούμενα : οὐδὲν γὰρ δίχα κινήσεως γίνεται. ξυνετέκνου δὲ ἀντὶ τοῦ ἐτεκνοποίει ; Hsch.~τ~384 τεκνοῦσα· τέκνον ἔμβρυον ἔχουσα}
\lp{16 ᾧ μὲν βλέπειν χρὴ πρῶτ᾽ ἐμηχανήσατο]}{Eur.~\textit{Supp.}~201–204 αἰνῶ δ᾽ ὃς ἡμῖν βίοτον ἐκ πεφυρμένου … πρῶτον μὲν ἐνθεὶς σύνεσιν, εἶτα δ᾽ ἄγγελον / γλῶσσαν λόγων δούς ; Parm.~28 B~13~DK (S.E. M.~9.9) πρώτιστον μὲν ἔρωτα θεῶν μητίσατο πάντων ; Emp.~31 B~86~DK ἐξ ὧν ὄμματ᾽ ἔπηξεν ἀτειρέα δῖ᾽ Ἀφροδίτη ; Emp.~31 B~87~DK γόμφοις ἀσκήσασα καταστόργοις Ἀφροδίτη ; Pl.~\textit{R.}~507c τὸν τῶν αἰσθήσεων δημιουργὸν ὅσῳ πολυτελεστάτην τὴν τοῦ ὁρᾶν τε καὶ ὁρᾶσθαι δύναμιν ἐδημιούργησεν ; Pl.~\textit{Ti.}~45b τῶν δὲ ὀργάνων πρῶτον μὲν φωσφόρα συνετεκτήναντο ὄμματα … σῶμα ἐμηχανήσαντο γίγνεσθαι ; X.~\textit{Mem.}~1.4.5 ὀφθαλμοὺς μὲν ὥσθ᾽ ὁρᾶν τὰ ὁρατά, ὦτα δὲ ὥστ᾽ ἀκούειν τὰ ἀκουστά ; Hsch.~μ~1301 μηχαναί· τέχναι. ἐπίνοιαι. βουλεύματα}
\lp{17 ὀφθαλμὸν ἀντίμιμον ἡλίου τροχῷ]}{Ar.~\textit{Nub.}~285–286 ὄμμα γὰρ αἰθέρος ἀκάματον σελαγεῖται / μαρμαρέαις ἐν αὐγαῖς ; Soph.~\textit{Ant.}~103 ἁμέρας βλέφαρον ; Emp.~31 B~84.7–8~DK ὣς δὲ τότ᾽ ἐν μήνιγξιν ἐεργμένον ὠγύγιον πῦρ / λεπτῇσίν ⟨τ᾽⟩ ὀθόνῃσι λοχάζετο κύκλοπα κούρην ; Pi. ap. D.H. Dem.~7 ἀκτὶς ἀελίου, τί πολύσκοπ᾽ ἐμήσω θοῶν μᾶτερ ὀμμάτων ; Pl.~\textit{R.}~508b ἀλλ᾽ ἡλιοειδέστατόν γε οἶμαι τῶν περὶ τὰς αἰσθήσεις ὀργάνων}
\lp{17 ἀντίμιμον]}{Sopat. ap.~Ath.~6.230f δρακοντομίμοις ἠρμένην τορεύμασιν ; Eur.~\textit{Ba.}~980 ἐπὶ τὸν ἐν γυναικομίμῳ στολᾷ ; Eur.~fr.~185.3~N.² (Antiope) γυναικομίμῳ διαπρέπεις μορφώματι ; Aesch.~\textit{PV}~1005 γυναικομίμοις ὑπτιάσμασιν χερῶν ; Alcid. ap. Arist. Rh.~1406a29 καὶ ἀντίμιμον τὴν τῆς ψυχῆς ἐπιθυμίαν ; Eratosth.~\textit{Cat.}~13 τῇ τοῦ Ἡλίου ἀντίμιμον ἐποιήσατο διφρείαν ; Ph.~\textit{Fug.}~179 μονονοὺκ ἀντίμιμον οὐρανοῦ δύναμιν ἐπιδείκνυσθαι δοκῶν ; Ph.~\textit{Somn.}~1.215 τοῦ παντὸς ἀντίμιμον ὄντα οὐρανοῦ ; Callistr.~\textit{Stat.}~9 εὐπαθοῦντι δὲ ἀνταποδιδόναι τὴν ἠχὴν ἀντίμιμον ; Hld.~9.9 ἀντίμιμον οὐρανοῦ τὸν ποταμὸν σεμνηγοροῦντες ; Σ \textit{Thesm.}~17 ἀντίμιμον : ὅμοιον. λείπει τὸ περιφερές ; Hsch.~α~5460 ἀντίπλαστον· Σοφοκλῆς Ἰνάχῳ … ἀντὶ τοῦ ἰσόπλαστον, ὅμοιον}
\lp{17 ἡλίου τροχῷ]}{Eur.~\textit{Phoen.}~3 Ἥλιε, θοαῖς ἵπποισιν εἱλίσσων φλόγα ; Aesch.~\textit{PV}~91 καὶ τὸν πανόπτην κύκλον ἡλίου καλῶ ; Soph.~\textit{Aj.}~845–846 σὺ δ᾽, ὦ τὸν αἰπὺν οὐρανὸν διφρηλατῶν / Ἥλιε ; Soph.~\textit{Ant.}~1064–1065 μὴ πολλοὺς ἔτι / τρόχους ἁμιλλητῆρας ἡλίου τελεῖν ; Anaximand. ap. [Plut.] Plac.~2.20,~889f ἁρματείῳ τροχῷ παραπλήσιον … καὶ τοῦτ᾽ εἶναι τὸν ἥλιον ; Hsch.~τ~1531 τροχόεν· τροχοειδές, στρογγύλον ; Hsch.~τ~1532 τροχός· περιβόλαιον, τεῖχος. ἢ κύκλος}
\lp{18 ἀκοὴν δὲ χοάνης ὦτα διετετρήνατο]}{Aesch.~\textit{Cho.}~451–452 δι᾽ ὤτων δὲ συν- / τέτραινε μῦθον ἡσύχῳ φρενῶν βάσει ; D.S.~2.56.4 καὶ τὰ μὲν τῆς ἀκοῆς τρήματα πολὺ τῶν παρ᾽ ἡμῖν ἔχειν εὐρυχωρέστερα ; Σ \textit{Thesm.}~18 λείπει ὡς}
\lp{18 χοάνης]}{\textit{Test.} Hsch.~χ~576 χοάνη· χώνη, τύπος, εἰς ὃν μεταχεῖται τὸ χωνευόμενον ; Moer. χ (p.~212 Bekker) χοάνην Ἀττικοί, χώνην Ἕλληνες ; Ar.~\textit{Thesm.}~57 καὶ χοανεύει. | καὶ λαικάζει ; Ar.~\textit{Thesm.}~62 τουτὶ τὸ πέος χοανεῦσαι ; Pherecr. fr.~113.30–31~K.-A. (108.30–31~Kock) πλήρεις κύλικας οἴνου μέλανος ἀνθοσμίου / ἤντλουν διὰ χώνης τοῖσι βουλομένοις πιεῖν ; Eup. fr.~283~Kock ἀλλ᾽ ὦ φίλε Ζεῦ κατάχυτλον τὴν ῥῖν᾽ ἔχεις ; Soph.~fr.~35~R. γλίδου χοάναι· εἰς οὓς τύπους κατέρχεται τὸ ἐν τῇ χωνῃ διὰ τῶν κεντημάτων. ἦσαν δὲ πήλινοι ; \textit{Il.}~18.470 φῦσαι δ᾽ ἐν χοάνοισιν ἐείκοσι πᾶσαι ἐφύσων ; Hes.~\textit{Th.}~863 τέχνῃ ὕπ᾽ αἰζηῶν ἐν ἐυτρήτοις χοάνοισι ; Emp.~31 B~84.9~DK ⟨αἳ⟩ χοάνῃσι δίαντα τετρήατο θεσπεσίῃσιν ; Emp.~31 B~96.1~DK (Arist. de An.~410a) ἡ δὲ χθὼν ἐπίηρος ἐν εὐστέρνοις χοάνοισιν ; Pl.~\textit{R.}~411a καταυλεῖν καὶ καταχεῖν τῆς ψυχῆς διὰ τῶν ὤτων ὥσπερ διὰ χώνης ; Ath.~10.436e ἐντιθέμενος γὰρ τῷ στόματι χώνην ἀπαύστως ἔπινεν ; Hsch.~ε~4248 ἐπάσκιον· ἡ χώνη. Σικελοί ; Hsch.~ε~7303 εὔχους· χώνη. Σαλαμίνιοι ; Hsch.~λ~960 λιγδαρεοχύται· οἱ ἐν τοῖς λίγδοις τὰς †σάρκας χέοντες, τουτέστιν χοάναις ; Hsch.~λ~965 λίγδος· τόπος χοάνης, καὶ ἡ θυεία, καὶ ἐν ᾗ χωνεύουσιν ; Hsch.~λ~966 λίγδου χοάνη· οἱ δὲ τὰ λίκνα τῶν ἀργυρείων ; Hsch.~ρ~342 ῥινοχόος· χώνη ; Hsch.~χ~575 χοανεῦσαι· χωνεῦσαι ; Hsch.~χ~577 χοάνοις· τοῖς φυσητῆρσι, ταῖς χώναις … ἢ τοῖς πηλίνοις τύποις ; Poll.~10.75 ἐκ δὲ τούτων καὶ ἡ χώνη, ὅταν εἴπῃ Φερεκράτης ἐν Πετάλῃ ; Poll.~10.147 χαλκέως δὲ σκεύη … χοάνη}
\lp{18 ἀκοὴν]}{Soph.~\textit{OT}~1386–1387 ἀλλ᾽ εἰ τῆς ἀκουούσης ἔτ᾽ ἦν / πηγῆς δι᾽ ὤτων φραγμός ; Pl.~\textit{Phdr.}~235d διὰ τῆς ἀκοῆς πεπληρῶσθαί με δίκην ἀγγείου ; Suid.~α~912 Ἀκοή: τὸ μέρος τοῦ σώματος}
\lp{18 ὦτα]}{\textit{Il.}~18.378–379 οὔατα δ᾽ οὔ πω / δαιδάλεα προσέκειτο ; Hsch.~ω~472 ὦτα· ἀκοαί, ὠτία ; Hsch.~ω~473 ὦτα· τὰ μὴ στρογγύλα}
\lp{18 διετετρήνατο]}{Emp.~31 B~100.3~DK (Arist. Resp.~473b) καί σφιν ἐπὶ στομίοις πυκναῖς τέτρηνται ἄλοξιν ; AP~4.3 (Agathias) εἰ μὴ φέροι πως ὦτα μὴ τετρημένα ; Hp.~\textit{Loc.\,Hom.}~2 Φύσις δὲ τοῦ σώματος, ἀρχὴ τοῦ ἐν ἰητρικῇ λόγου· πρῶτον διατέτρηται ᾗ ἐσακούομεν ; Thphr.~\textit{Sens.}~41 (Diog. Apoll.) ἡ ἀκοὴ τέτρηται βραχὺ καὶ λεπτὸν καὶ ἰθύ ; Plu.~Mor.~502d μήποτε γὰρ αὐτοῖς οὐκ εἰς τὴν ψυχὴν ἀλλ᾽ εἰς τὴν γλῶτταν ἡ ἀκοὴ συντέτρηται ; Arr. Epict.~2.24.12 ἔχομεν καὶ κατὰ ποσὸν περὶ τοῦτο τὰ ὦτα τετρημένα ; Hsch.~τ~645 τετρῆναι· τρῆσαι ; Hsch.~τ~1315 τρήματα· τρυπήματα}
\lp{19 διὰ τὴν χοάνην οὖν]}{Ar.~\textit{Nub.}~227 ἀλλ᾽ οὐκ ἀπὸ τῆς γῆς, εἴπερ; | οὐ γὰρ ἄν ποτε ; Ar.~\textit{Nub.}~236 ἡ φροντὶς ἕλκει τὴν ἰκμάδ᾽ ἐς τὰ κάρδαμα ; Ar.~\textit{Nub.}~380–381 τουτί μ᾽ ἐλελήθειν, / ὁ Ζεὺς οὐκ ὤν, ἀλλ᾽ ἀντ᾽ αὐτοῦ Δῖνος νυνὶ βασιλεύων}
\lp{19 μήτ᾽ ἀκούω μήθ᾽ ὁρῶ]}{Parm.~28 B~6.7~DK οἱ δὲ φοροῦνται κωφοὶ ὁμῶς τυφλοί τε, τεθηπότες, ἄκριτα φῦλα ; Tosi~1971 Nec aures me credo habere nec tango}
\lp{20 νὴ τὸν Δί᾽ ἥδομαί γε]}{Ar.~\textit{Nub.}~153 ὦ Ζεῦ βασιλεῦ τῆς λεπτότητος τῶν φρενῶν ; Ar.~\textit{Nub.}~773 οἴμ᾽ ὡς ἥδομαι ; Ar.~\textit{Lys.}~1246 ὡς ἥδομαί γ᾽ ὑμᾶς ὁρῶν ὀρχουμένους}
\lp{20 προσμαθών]}{Ar.~\textit{Thesm.}~24 ἔτι προσμάθοιμι χωλὸς εἶναι τὼ σκέλει ; Aesch.~\textit{PV}~697 ἐπίσχες ἔστ᾽ ἂν καὶ τὰ λοιπὰ προσμάθῃς ; Aesch.~fr.~396~R. Καλὸν δὲ καὶ γέροντα μανθάνειν σοφά ; Sol. ap. Pl. Amat.~133c γηράσκω δ᾽ αἰεὶ πολλὰ διδασκόμενος ; Plu.~\textit{Sol.}~31.3 γηράσκω δ᾽ αἰεὶ πολλὰ διδασκόμενος ; Zen.~3.4 Γηράσκω δ᾽ αἰεὶ πολλὰ διδασκόμενος: ἐπὶ τῶν διὰ τὸ γῆρας ἐμπειροτέρων ; Diogenian.~3.80 Γηράσκω αἰεὶ πολλὰ διδασκόμενος: ἐπὶ τῶν διὰ γῆρας ἐμπειροτέρων ; Suid.~γ~246 καὶ παροιμία· Γηράσκω δ᾽ αἰεὶ πολλὰ διδασκόμενος, ἐπὶ τῶν διὰ γῆρας ἐμπειροτέρων ; Tosi~2035 Γηράσκω δ᾽ αἰεὶ πολλὰ διδασκόμενος}
\lp{21 οἷόν γέ που ᾿στιν]}{Ar.~\textit{Vesp.}~27 δεινόν γέ ποὔστ᾽ ἄνθρωπος ἀποβαλὼν ὅπλα ; Ar.~\textit{Pax}~828 οὔκ, εἰ μή γέ που ; Ar.~\textit{Pax}~1047 οὗτός γέ πού ᾽σθ᾽ ὁ χρησμολόγος οὑξ Ὠρεοῦ ; Ar.~\textit{Eccl.}~1035 νὴ τὸν Δί᾽ ἤνπερ ᾖ γέ που τῶν κηρίνων ; Eup. fr.~342~K.-A. (314~Kock) οἷόν γέ πού ᾽στι γλῶσσα κἀνθρώπου λόγος}
\lp{21 αἱ σοφαὶ ξυνουσίαι]}{Ar.~fr.~323~K.-A. (308~Kock) ἐν δὲ τοῖς Ἥρωσιν ἄντικρυς ἀποφαίνεται ; Soph.~fr.~14~R. (13~N.²) σοφοὶ τύραννοι τῶν σοφῶν ξυνουσίᾳ ; Pl.~\textit{R.}~568b ὡς ἄρα σοφοὶ τύραννοί εἰσι τῶν σοφῶν συνουσίᾳ ; Pl.~\textit{Thg.}~125b Εὐριπίδης γάρ πού φησιν— σοφοὶ τύραννοι τῶν σοφῶν συνουσίᾳ ; Aristid.~\textit{Or.}~46 (II p.~288 Dind.) ἔστι γὰρ ἐξ Αἴαντος Σοφοκλέους, Αἴαντος τοῦ Λοκροῦ ; Them.~\textit{Or.}~6, p.~72c Εὐριπίδης μὲν γὰρ ἢ ὅστις δή ποτέ ἐστιν ὁ ποιήσας σοφοὶ τύραννοι τῶν σοφῶν συνουσίᾳ ; Lib.~\textit{Ep.}~35.3 τὸ τοῦ Σοφοκλέους σοφοὶ τύραννοι τῇ σοφῶν συνουσίᾳ ; Stob.~\textit{Flor.}~48.5~Meineke Εὐριπίδου. Σοφοὶ τύραννοι τῶν σοφῶν συνουσίᾳ ; Gell.~\textit{NA}~13.19.1–2 Eum versum Plato in Theaeteto Euripidi esse dicit ; Σ \textit{Thesm.}~21 διὰ τούτου ὑπονοεῖ Εὐριπίδου εἶναι ἐκεῖνο τὸ « σοφοὶ τύραννοι τῶν σοφῶν ξυνουσίᾳ. » ἔστι δὲ Σοφοκλέους ἐξ Αἴαντος Λοκροῦ … ἐν δὲ τοῖς Ἥρωσιν ἄντικρυς ἀποφαίνεται ; Zen.~5.98 Σοφοὶ τύραννοι τῶν σοφῶν συνουσίᾳ: τοῦτο Σοφοκλέους ἐστὶν ἐξ Αἴαντος τοῦ Λοκροῦ· Πλάτων δὲ Εὐριπίδου φησὶν εἶναι τὸ ἰαμβεῖον ; App. Prov.~4.76 Σοφοὶ τύραννοι τῶν σοφῶν συνουσίᾳ ; Suid.~σ~823 Σοφοὶ τύραννοι τῇ τῶν σοφῶν συνουσίᾳ}
\lp{21 αἱ σοφαί]}{Hsch.~σ~1371 σοφιστήν· πᾶσαν τέχνην σοφίαν ἔλεγον ; Hsch.~σ~1373 σοφός· φρόνιμος. φιλότεχνος, καὶ ἐξευρετικός. καὶ ὁ τῶν θείων ἔμπειρος}
\lp{21 ξυνουσίαι]}{Ar.~\textit{Nub.}~649 πρῶτον μὲν εἶναι κομψὸν ἐν συνουσίᾳ ; Ar.~\textit{Vesp.}~1256 οὔκ, ἢν ξυνῇς γ᾽ ἀνδράσι καλοῖς τε κἀγαθοῖς ; Ar.~\textit{Av.}~324 ἄνδρ᾽ ἐδεξάμην ἐραστὰ τῆσδε τῆς ξυνουσίας ; Ar.~\textit{Eccl.}~110 καὶ πῶς γυναικῶν θηλύφρων ξυνουσία ; Thgn. ap. X. Mem.~1.2.20 ἐσθλῶν μὲν γὰρ ἄπ᾽ ἐσθλὰ διδάξεαι· ἢν δὲ κακοῖσι / συμμίσγῃς, ἀπολεῖς καὶ τὸν ἐόντα νόον ; Pl.~\textit{Prt.}~316c ἀπολείποντας τὰς τῶν ἄλλων συνουσίας … ἑαυτῷ συνεῖναι ; Pl.~\textit{Thg.}~122e ἀνδρός τινος συνουσίαν τοιούτου ὅστις σε σοφὸν ποιήσει ; X.~\textit{Mem.}~4.2.2 πότερον Θεμιστοκλῆς διὰ συνουσίαν τινὸς τῶν σοφῶν ἢ φύσει τοσοῦτον διήνεγκε τῶν πολιτῶν ; 1 \textit{Ep.\,Cor.}~15.33 φθείρουσιν ἤθη χρηστὰ ὁμιλίαι κακαί ; Hsch.~ξ~168 ξυνουσία· συνήθεια. συνουσία ; Suid.~ξ~148 Ξυνουσίᾳ: τῷ ἔργῳ. τῇ πείρᾳ}
\lp{22 πόλλ᾽ ἂν μάθοις … παρ᾽ ἐμοῦ]}{Ar.~\textit{Nub.}~459–460 ταῦτα μαθὼν παρ᾽ ἐμοῦ κλέος οὐρανόμηκες / ἐν βροτοῖσιν ἕξεις}
\lp{22 πῶς ἂν οὖν]}{Ar.~\textit{Eq.}~17–18 πῶς ἂν οὖν ποτε / εἴποιμ᾽ ἂν αὐτὸ δῆτα κομψευριπικῶς ; Ar.~\textit{Pax}~68 πῶς ἄν ποτ᾽ ἀφικοίμην ἂν εὐθὺ τοῦ Διός ; Ar.~\textit{Thesm.}~1015–1018 φίλαι παρθένοι φίλαι, πῶς ἂν ⟨οὖν⟩ / … / τὸν Σκύθην λάθοιμι ; Eur.~\textit{Alc.}~864 πῶς ἂν ὀλοίμαν ; Eur.~\textit{Med.}~97 ἰώ μοί μοι, πῶς ἂν ὀλοίμαν}
\lp{[23] πρὸς τοῖς ἀγαθοῖς τούτοισιν]}{Ar.~\textit{Pax}~369–370 κᾆτα τῷ τρόπῳ / οὐκ ᾐσθόμην ἀγαθὸν τοσουτονὶ λαβών ; Ar.~\textit{Plut.}~462 τί δ᾽ ἂν ὑμεῖς ἀγαθὸν ἐξεύροιθ᾽ ; Soph.~\textit{Ant.}~31 τοιαῦτά φασι τὸν ἀγαθὸν Κρέοντα σοὶ / κἀμοί … ; Soph.~\textit{Ant.}~274–275 κἀμὲ τὸν δυσδαίμονα / πάλος καθαιρεῖ τοῦτο τἀγαθὸν λαβεῖν ; Soph.~\textit{Phil.}~872–873 οὔκουν Ἀτρεῖδαι τοῦτ᾽ ἔτλησαν εὐφόρως / οὕτως ἐνεγκεῖν, ἁγαθοὶ στρατηλάται ; Σ Soph.~\textit{Ant.}~31 ἀγαθόν: ἐν εἰρωνείᾳ}
\lp{[23]–24 ἐξεύροιμ᾽ ὅπως … προσμάθοιμι]}{Ar.~\textit{Pax}~31–32 μὴ παύσαιο μηδέποτ᾽ ἐσθίων / τέως ἕως σαυτὸν λάθοις διαρραγείς ; Ar.~\textit{Pax}~521–522 πόθεν ἂν λάβοιμι ῥῆμα μυριάμφορον / ὅτῳ προσείπω σ᾽ ; Ar.~\textit{Ran.}~96–97 γόνιμον δὲ ποιητὴν ἂν οὐχ εὕροις ἔτι / ζητῶν ἄν, ὅστις ῥῆμα γενναῖον λάκοι ; Soph.~\textit{Aj.}~1217–1222 γενοίμαν ἵν᾽ ὑλᾶεν ἔπεστι πόντου / πρόβλημ᾽ … / τὰς ἱερὰς ὅπως / προσείποιμεν Ἀθάνας}
\lp{24 χωλὸς εἶναι]}{Ar.~\textit{Ach.}~410–411 ἀναβάδην ποιεῖς, / ἐξὸν καταβάδην; οὐκ ἐτὸς χωλοὺς ποιεῖς ; Ar.~\textit{Pax}~147–148 εἶτα χωλὸς ὢν Εὐριπίδῃ / λόγον παράσχῃς καὶ τραγῳδία γένῃ ; Ar.~\textit{Av.}~1292–1293 πέρδιξ μὲν εἷς κάπηλος ὠνομάζετο / χωλός ; Ar.~\textit{Ran.}~846 τὸν χωλοποιὸν οἷος ὢν θρασύνεται ; Ar.~fr.~53~Kock (Anagyros) καὶ μὴν χθές γ᾽ ἦν Πέρδιξ χωλός ; Timocl. fr.~6.15~Kock (Dionysiazousai) χωλός τίς ἐστι, τὸν Φιλοκτήτην ὁρᾷ ; Σ \textit{Thesm.}~24 ἵνα μὴ περιπατήσῃ μετ᾽ αὐτοῦ ; Hsch.~χ~858 χωλεύει· νοσεῖ. σκάζει ; App. Prov.~2.65 Ἐν παντὶ μύθῳ καὶ τὸ Πέρδικος σκέλος· … Πέρδιξ γὰρ ἦν τις Ἀθήνησι χωλὸς κάπηλος}
\lp{24 χωλὸς εἶναι (proverbe)]}{Com. adesp. fr.~610~Kock χωλῷ παροικῶν κἂν ἑνὶ σκάζειν μάθοις ; Aesch.~\textit{Pers.}~753–754 ταῦτά τοι κακοῖς ὁμιλῶν ἀνδράσιν διδάσκεται / θούριος Ξέρξης ; Aesch.~\textit{Sept.}~599–600 ἐν παντὶ πράγει δ᾽ ἔσθ᾽ ὁμιλίας κακῆς / κάκιον οὐδέν ; Thgn.~35–36 ἐσθλῶν μὲν γὰρ ἄπ᾽ ἐσθλὰ μαθήσεαι· ἢν δὲ κακοῖσιν / συμμίσγῃς, ἀπολεῖς καὶ τὸν ἐόντα νόον ; Plu.~De lib. educ.~6 (Mor.~4a) καὶ οἱ παροιμιαζόμενοι δέ φασιν … ὅτι ἂν χωλῷ παροικήσῃς, ὑποσκάζειν μαθήσῃ ; Σ Pind.~\textit{N.}~7.127c ὥσπερ ἡ παροιμία ἔχει· χωλῷ παροικεῖς, κἂν ἑνὶ σκάζειν μάθοις ; Macar.~6.90 Παρὰ χωλὸν οἰκῶν κἂν ἐπισκάζειν μαθήσῃ· πρὸς τοὺς παραπολαύοντας τῶν χαλεπῶν ; Apostol.~2.94 Ἂν χωλῷ παροικήσῃς, ὑποσκάζειν μαθήσῃ· ὅτι οἱ βαρβάροις καὶ μοχθηροῖς συναναχρωννύμενοι ἀποφέρονταί τι τῆς ἐκείνων φαυλότητος ; App. Prov.~1.60 Βότρυς πρὸς βότρυν πεπαίνεται· ἐπὶ τῶν ἐξισοῦσθαι φιλονεικούντων ; Tosi~461 Φθείρουσιν ἤθη χρῆσθ᾽ ὁμιλίαι κακαί ; Tosi~1647 Παρὰ χωλὸν οἰκῶν κἂν ἐπισκάζειν μαθήσῃ}
\lp{24 τὼ σκέλει]}{Ar.~\textit{Pax}~820 ἔγωγέ τοι πεπόνηκα κομιδῇ τὼ σκέλει ; Ar.~\textit{Pax}~825–826 ἤλγουν τὼ σκέλει μακρὰν ὁδὸν / διεληλυθώς ; Alex. fr.~147.3~Kock σοφὸν οὐδὲν εὕρηκ᾽, ἀλλὰ κοπιῶ τὰ σκέλη ; Pl.~\textit{Smp.}~190d ἐφ᾽ ἑνὸς πορεύσονται σκέλους ἀσκωλιάζοντες ; Hsch.~σ~909 σκελοῖν· ποδῶν ; Hsch.~σ~910 σκέλος· τάττεται ἐπὶ τοῦ φορτικοῦ. καὶ μέρος τι τῆς νεώς}
\lp{25 βάδιζε δευρί]}{Hsch.~β~34 βαδίζει· περιπατεῖ}
\lp{25 πρόσεχε τὸν νοῦν — ἰδού]}{Ar.~\textit{Eq.}~1014 ἄκουε δή νυν καὶ πρόσεχε τὸν νοῦν ἐμοί ; Ar.~\textit{Nub.}~254–255 κάθιζε τοίνυν ἐπὶ τὸν ἱερὸν σκίμποδα. — ἰδοὺ κάθημαι ; Ar.~\textit{Nub.}~635 ἀνύσας τι κατάθου καὶ πρόσεχε τὸν νοῦν. — ἰδού ; Ar.~\textit{Thesm.}~381 σίγα σιώπα, πρόσεχε τὸν νοῦν· χρέμπτεται γὰρ ἤδη ; Σ \textit{Thesm.}~25 πρόσεχε ἀντὶ τοῦ μνῆσαι ; Hsch.~π~3753 προσεχές· ἐγγύς. νῆφον, νουνεχές}
\lp{25–28 πρόσεχε τὸν νοῦν … σίγα … ἄκου᾽]}{Cratin. fr.~315~K.-A. (=~284~Kock) ἄκουε, σίγα, πρόσεχε τὸν νοῦν, δεῦρ᾽ ὅρα ; Soph.~fr.~815 Radt (=~1014~N.²) ἄκουε, σίγα· τίς ποτ᾽ ἐν δόμοις βοή ; Synag. Σᵇ α~762~Cunningham ἄκουε· ἀντὶ τοῦ σίγα. Κρατῖνος (fr. 315). οἷον τόπον δίδου διὰ τῆς σιγῆς τῇ ἀκοῇ}
\lp{26 ὁρᾷς τὸ θύριον τοῦτο]}{Ar.~\textit{Nub.}~91–92 δεῦρό νυν ἀπόβλεπε. / ὁρᾷς τὸ θύριον τοῦτο καὶ τᾠκίδιον}
\lp{26 νὴ τὸν Ἡρακλέα]}{Ar.~\textit{Av.}~1391 οὐ δῆτ᾽ ἔγωγε. — νὴ τὸν Ἡρακλέα σύ γε}
\lp{27–28 σιωπῶ τὸ θύριον]}{Ar.~\textit{Eccl.}~549–550 μὴ φροντίσῃς· ἄρρεν γὰρ ἔτεκε παιδίον. — ἡκκλησία; — μὰ Δί᾽ ἀλλ᾽ ἐφ᾽ ἣν ἐγᾠχόμην ; Ar.~\textit{Plut.}~1098–1099 ἀλλὰ δῆτα τὸ θύριον / φθεγγόμενον ἄλλως κλαυσιᾷ ; Thgn.~421–422 πολλοῖς ἀνθρώπων γλώσσῃ θύραι οὐκ ἐπίκεινται / ἁρμόδιαι, καί σφιν πόλλ᾽ ἀμέλητα μέλει ; X.~\textit{Smp.}~6.10 καὶ ῥᾳδίως γ᾽, ἂν ἃ μὴ δεῖ λέγειν, ἔφη, σιωπᾷς ; Plu.~De fac.~26,~940f ἐπίσχες εἶπεν ὦ Λαμπρία, καὶ παραβαλοῦ τὸ θυρίον τοῦ λόγου ; Σ \textit{Thesm.}~27 λείπει ἡ διά, οἷον σιωπῶ διὰ τὸ θύριον. ἔπαιξε δέ ; Hsch.~σ~473 σιγᾷ· σιωπᾷ}
\lp{29 ἐνταῦθ᾽ … οἰκῶν τυγχάνει]}{Ar.~\textit{Nub.}~94–95 ψυχῶν σοφῶν τοῦτ᾽ ἐστὶ φροντιστήριον. / ἐνταῦθ᾽ ἐνοικοῦσ᾽ ἄνδρες …}
\lp{29 Ἀγάθων]}{Ar.~\textit{Ran.}~83–84 Ἀγάθων δὲ ποῦ ᾽στιν; — ἀπολιπών μ᾽ ἀποίχεται, / ἀγαθὸς ποιητὴς καὶ ποθεινὸς τοῖς φίλοις ; Ar.~fr.~341~K.-A. (=~326~Kock) καὶ κατ᾽ Ἀγάθων᾽ ἀντίθετον ἐξυρημένον ; Pl.~\textit{Smp.}~173a ὅτε τῇ πρώτῃ τραγῳδίᾳ ἐνίκησεν Ἀγάθων ; Arist.~\textit{Po.}~1456a29–30 διὸ ἐμβόλιμα ᾄδουσιν πρώτου ἄρξαντος Ἀγάθωνος τοῦ τοιούτου ; Luc.~\textit{Rh.\,Pr.}~11 αὐτὸν Ἀγάθωνα, τὸν τῆς τραγῳδίας ἐπέραστον ἐκεῖνον ποιητήν ; Ael.~\textit{VH}~14.13 ἀλλὰ σύ γε, γενναῖε, λέληθας σεαυτὸν τὸν Ἀγάθωνα ἐκ τοῦ Ἀγάθωνος ἀφανίζων ; Ath.~5.217a ὅτε γὰρ Ἀγάθων ἐνίκα … ἐπὶ ἄρχοντος Εὐφήμου στεφανοῦται Ληναίοις ; Σ Ar.~\textit{Ran.}~83–84 κωμῳδεῖται δὲ εἰς θηλύτητα … οὗτος δὲ ἀγαθὸς ἦν τὸν τρόπον καὶ τὴν τράπεζαν λαμπρός ; Σ Luc.~\textit{Rh.\,Pr.}~11 Ἀγάθων τραγῳδίας ποιητὴς εἰς μαλακίαν σκωπτόμενος … παιδικὰ γεγονὼς Παυσανίου τοῦ τραγικοῦ ; Suid.~α~124 τραγικὸς δὲ ἦν· διεβέβλητο δὲ ἐπὶ μαλακίᾳ … οὗτος ὁ Ἀγάθων ἀγαθὸς ἦν τὸν τρόπον}
\lp{29 ὁ κλεινός]}{Ar.~\textit{Ach.}~1184 ὦ κλεινὸν ὄμμα νῦν πανύστατόν σ᾽ ἰδὼν ; Ar.~\textit{Eq.}~1328 ἵν᾽ ὁ κλεινὸς Δῆμος ἐνοικεῖ ; Ar.~\textit{Plut.}~772 ἔπειτα σεμνῆς Παλλάδος κλεινὸν πέδον ; Eur.~\textit{El.}~327 μέθῃ δὲ βρεχθεὶς τῆς ἐμῆς μητρὸς πόσις / ὁ κλεινός, ὡς λέγουσιν, ἐνθρῴσκει τάφῳ ; Eur.~\textit{Or.}~17 ὁ κλεινός, εἰ δὴ κλεινός, Ἀγαμέμνων ἔφυ ; Soph.~\textit{El.}~300–301 ὁ κλεινὸς αὐτῇ ταὐτὰ νυμφίος παρών, / ὁ πάντ᾽ ἄναλκις οὗτος ; Pl.~\textit{Smp.}~175e ἐν μάρτυσι τῶν Ἑλλήνων πλέον ἢ τρισμυρίοις — ὑβριστὴς εἶ, ἔφη, ὦ Σώκρατες, ὁ Ἀγάθων ; Σ Soph.~\textit{El.}~300 ὁ [κλεινὸς] κατὰ εἰρωνείαν ; Hsch.~κ~2900 κλεινός· ὡραῖος, ἔνδοξος, ὀνομαστός ; Hsch.~κ~2902 κλεινοί· οἱ εἰς τὰ παιδικὰ ἐπὶ κάλλει ἁρπαζόμενοι παῖδες ; Hsch.~κ~3006 κλεινός· ἔνδοξος, ἀγαθός, καλός ; Suid.~κ~1752 Κλεινός: ἔνδοξος· ἀπὸ τοῦ κλέους}
\lp{30 ὁ τραγῳδοποιός]}{Ar.~\textit{Thesm.}~88 Ἀγάθωνα πεῖσαι τὸν τραγῳδοδιδάσκαλον ; Pl.~\textit{Smp.}~223d καὶ τὸν τέχνῃ τραγῳδοποιὸν ὄντα ⟨καὶ⟩ κωμῳδοποιὸν εἶναι}
\lp{30 ποῖος οὗτος Ἁγάθων]}{Ar.~\textit{Ach.}~963–964 ὁ ποῖος οὗτος Λάμαχος τὴν ἔγχελυν; — ὁ δεινός, ὁ ταλαύρινος, ὃς τὴν Γοργόνα … ; Ar.~\textit{Thesm.}~96 καὶ ποῖός ἐστιν; — οὗτος οὑκκυκλούμενος ; Ar.~\textit{Thesm.}~620–621 τὸν δεῖνα γιγνώσκεις, τὸν ἐκ Κοθωκιδῶν; — τὸν δεῖνα; ποῖον ; Ar.~\textit{Plut.}~392 σὺ Πλοῦτον; ποῖον; — αὐτὸν τὸν θεόν ; Pherecr. fr.~155.20–21~K.-A. (145~Kock) ποῖος οὑτοσὶ / ὁ Τιμόθεος; — Μιλήσιός τις πυρρίας ; Aeschin.~1.130 οὐχ ἅμα τοὔνομα λέγεται καὶ τὸ ἐρώτημα ἐρωτᾶτε· ποῖος Τίμαρχος; ὁ πόρνος ; Charito~5.10.1 ποῖος οὗτος ἐπ᾽ ἐμοῦ Πρωτεσίλεως ἀνεβίω}
\lp{31 ἔστιν τις ἀγαθῶν]}{Ar.~\textit{Vesp.}~181 Ὀδυσσέα τιν᾽ ; Pl.~\textit{Prt.}~315d–e καλόν τε κἀγαθὸν τὴν φύσιν … ὄνομα αὐτῷ εἶναι Ἀγάθωνα, καὶ οὐκ ἂν θαυμάζοιμι εἰ παιδικὰ Παυσανίου τυγχάνει ὤν ; Pl.~\textit{Smp.}~174b ἵνα καὶ τὴν παροιμίαν διαφθείρωμεν μεταβαλόντες, ὡς ἄρα καὶ Ἀγάθων᾽ ἐπὶ δαῖτας ἴασιν αὐτόματοι ἀγαθοί ; Aeschin.~1.41 Μισγόλας ἔστι τις … ἀνὴρ τὰ μὲν ἄλλα καλὸς κἀγαθός ; Aeschin.~1.62 ἔστι τις Γλαύκων Χολαργεύς ; Hsch.~α~277 ἀγαθόν· ἀνδρεῖον, γενναῖον ; Hsch.~ν~152 πολλὰ δέ ἐστι τοιαῦτα· καὶ ἀγαθῶν ⟨Ἀγαθίδες⟩ καὶ λευκότερος Λεύκωνος ; Zen.~2.19 Αὐτόματοι δ᾽ ἀγαθοὶ ἀγαθῶν ἐπὶ δαῖτας ἴενται … Καὶ ὁ Πλάτων ἐν τῷ Συμποσίῳ οὕτως αὐτῇ ἐχρήσατο ; Suid.~α~123 Ἀγαθῶν ἀγαθίδες: τάττεται ἡ παροιμία παρὰ τοῖς κωμικοῖς ἐπὶ τῶν πολλῶν ἀγαθῶν ; Tosi~516 Αὐτόματοι δ᾽ ἀγαθοὶ ἀγαθῶν ἐπὶ δαῖτας ἴενται}
\lp{31 μῶν ὁ μέλας ὁ καρτερός]}{Ar.~\textit{Lys.}~800–804 καὶ Μυρωνίδης γὰρ ἦν / τραχὺς ἐντεῦθεν μελάμπυγός / τε τοῖς ἐχθροῖς ἅπασιν, / ὣς δὲ καὶ Φορμίων ; Ar.~\textit{Thesm.}~191–192 σὺ δ᾽ εὐπρόσωπος λευκὸς ἐξυρημένος / γυναικόφωνος ἁπαλὸς εὐπρεπὴς ἰδεῖν ; Ar.~\textit{Ran.}~1091–1092 βραδὺς ἄνθρωπός τις ἔθει κύψας / λευκὸς πίων ὑπολειπόμενος ; Ar.~\textit{Eccl.}~60–67 ἀλειψαμένη τὸ σῶμ᾽ ὅλον δι᾽ ἡμέρας / ἐχραινόμην ἑστῶσα πρὸς τὸν ἥλιον … τὸ ξυρὸν δέ γ᾽ ἐκ τῆς οἰκίας / ἔρριψα πρῶτον, ἵνα δασυνθείην ὅλη ; Ar.~\textit{Eccl.}~427–428 μετὰ τοῦτο τοίνυν εὐπρεπὴς νεανίας / λευκός τις ἀνεπήδησ᾽ ὅμοιος Νικίᾳ ; Eur.~\textit{Ba.}~457–458 λευκὴν δὲ χροιὰν ἐκ παρασκευῆς ἔχεις, / οὐχ ἡλίου βολαῖσιν, ἀλλ᾽ ὑπὸ σκιᾶς ; \textit{Od.}~16.175–176 ἂψ δὲ μελαγχροιὴς γένετο, γναθμοὶ δὲ τάνυσθεν, / κυάνεαι δ᾽ ἐγένοντο γενειάδες ἀμφὶ γένειον ; Pl.~\textit{Phdr.}~239c–d μαλθακόν τινα καὶ οὐ στερεὸν διώκων, οὐδ᾽ ἐν ἡλίῳ καθαρῷ τεθραμμένον ἀλλὰ ὑπὸ συμμιγεῖ σκιᾷ ; Pl.~\textit{R.}~474e μέλανας δὲ ἀνδρικοὺς ἰδεῖν, λευκοὺς δὲ θεῶν παῖδας εἶναι ; Pl.~\textit{R.}~556d ἰσχνὸς ἀνὴρ πένης, ἡλιωμένος, παραταχθεὶς ἐν μάχῃ πλουσίῳ ἐσκιατροφηκότι ; X.~\textit{HG}~3.4.19 λευκοὺς μὲν διὰ τὸ μηδέποτε ἐκδύεσθαι, μαλακοὺς δὲ καὶ ἀπόνους ; D.~21.71 Σώφιλον τὸν παγκρατιαστήν (ἰσχυρός τις ἦν, μέλας, εὖ οἶδ᾽ ὅτι γιγνώσκουσίν τινες ὑμῶν ὃν λέγω) ; [Arist.] \textit{Phgn.}~6.32 Οἱ δὲ λευκοὶ ἄγαν δειλοί· ἀναφέρεται ἐπὶ τὰς γυναῖκας ; Σ \textit{Thesm.}~31 οἱ περὶ Ἀρίσταρχον καὶ Δίδυμόν φασιν εἶναι τοιοῦτον Ἀγάθωνα … ἐπειδὴ βούλεται κωμῳδεῖν τὸν Ἀγάθωνα ὡς μήτε καρτερὸν μήτε μέλανα ; Σ \textit{Il.}~4.141 ἐπὶ παντὸς γὰρ ζῴου τὸ μελάγχροόν ἐστιν ἰσχυρότερον ; Σ \textit{Od.}~16.175 στεριφώτερα δὲ τὰ μέλανα τῶν σωμάτων ; Hsch.~κ~899 καρτερός· κραταιός, ἰσχυρός. ἐγκρατής. κρατῶν, καὶ ἰσχύων ; Suid.~λ~713 ἐχλιαίνετο πρὸς τὸν ἥλιον ὥστε μελανθῆναι ὡς ἀνήρ ; App. Prov.~3.61 Λευκοὶ τοὺς δειλοὺς ὀνειδίζουσιν … παρόσον οἱ λευκοὶ δειλοί}
\lp{31 ὁ μέλας (μελάμπυγος)]}{Eub. fr.~61~Kock ἐγὼ δέ γ᾽ εἰμὶ τῶν μελαμπύγων ἔτι ; Alex. fr.~322~K.-A. λευκόπυγος· ὁ ἄνανδρος ; Archil. fr.~178 W. μή τευ μελαμπύγου τύχῃς ; Luc.~\textit{Pseudol.}~32 εἴ τις ἀνὴρ δασὺς καὶ τοῦτο δὴ τὸ ἀρχαῖον, μελάμπυγος ; Philostr.~\textit{VA}~2.36 τοῦτ᾽ ἴσως ἦν τὸ μελαμπύγου τυχεῖν ; Σ Ar.~\textit{Lys.}~802,~804 τοὺς λευκοπύγους ὡς γυναικώδεις ἐκωμῴδουν ; Eust.~863,~29 Ἄλεξις ὁ κωμικὸς λευκόπυγον ἔφη τὸν ἄνανδρον, οὗ ἔμπαλιν μελαμπύγους τοὺς ἀνδρείους ἔλεγον, ὅθεν καὶ παροιμία τό, οὔπω μελαμπύγῳ ἐνέτυχες ; Hsch.~λ~739 λευκόπυγος· ὁ ἄνανδρος. ἔμπαλιν δὲ μελαμπύγους τοὺς ἀνδρείους ἔλεγον ; Hsch.~μ~641 μελάμπυγος· ἀνδρεῖος. τοὺς γὰρ δασεῖς τὰς πυγὰς ἀνδρείους ἐνόμιζον ; Hsch.~μ~1277 μή τευ μελαμπύγου τύχοις· μή τινος ἀνδρείου καὶ ἰσχυροῦ τύχοις ; Suid.~λ~335 Λευκοπύγους: δειλούς, ὡς μελαμπύγους τοὺς ἀνδρείους ; Suid.~λ~715 οὗτος δὲ ὁ Φορμίων τριχώδης ἦν καὶ πυγὰς μελαίνας εἶχε ; Suid.~μ~449 Μελαμπύγου τύχοις: … ἡ μήτηρ ἔλεγε φυλάσσεσθαι μήποτε ἐμπέσωσιν εἰς δασύπρωκτον ; Suet.~\textit{Blasph.}~53~Taillardat Λευκόπυγος· ὁ ἄνανδρος· καὶ ἔμπαλιν μελάμπυγος, ὁ ἀνδρεῖος ; Zen.~5.10 Μὴ σύ γε μελαμπύγου τύχοις: μή τινος ἀνδρείου καὶ ἰσχυροῦ τύχοις ; Diogenian.~6.38 Μὴ μελαμπύγου τύχῃς: μή τινος ἀνδρείου καὶ ἰσχυροῦ ; App. Prov.~3.62 Λευκόπυγος: ἐπὶ τῶν δειλῶν καὶ ἀνάνδρων … Τούτῳ δὲ ἐναντίον τὸ μελάμπυγος, ἐπὶ τῶν γενναίων}
\lp{32 οὔκ, ἀλλ᾽ ἕτερός τις]}{Eup. fr.~182~Kock (Marikas) ἄκουε νῦν Πείσανδρος ὡς ἀπόλλυται. / Β. ὁ στρεβλός; Α. οὐκ ἀλλ᾽ ὁ μέγας, οὑνοκίνδιος ; Luc.~\textit{DMeretr.}~10.1 Οὔκ, ἀλλ᾽ ὁ κάκιστα φιλοσόφων ἀπολούμενος Ἀρισταίνετος. — Τὸν σκυθρωπὸν λέγεις, τὸν δασύν, τὸν βαθυπώγωνα ; Diogenian.~1.63 Ἄλλος οὗτος Ἡρακλῆς: ἐπὶ τῶν ἰσχυρῶν καὶ κραταιῶν ἡ παροιμία}
\lp{32 οὐχ ἑόρακας πώποτε]}{Ar.~\textit{Nub.}~767–768 ἤδη παρὰ τοῖσι φαρμακοπώλαις τὴν λίθον / ταύτην ἑόρακας τὴν καλήν, τὴν διαφανῆ ; Ar.~\textit{Av.}~1439 πάντες; — οὐκ ἀκήκοας ; Ar.~\textit{Ran.}~1323–1324 ὁρᾷς τὸν πόδα τοῦτον; — ὁρῶ. — τί δαί; τοῦτον ὁρᾷς; — ὁρῶ ; Σ \textit{Thesm.}~32 ἐπειδὴ οὐ πάλαι ἤρξατο διδάσκειν, ἀλλὰ τρισὶν πρὸ τούτων ἔτεσιν}
\lp{33 μῶν ὁ δασυπώγων]}{\textit{Test.} Poll.~2.88 παρὰ δὲ τοῖς ποιηταῖς, εὐπώγων τις εἴρηται, ἀλλ᾽ ἔστι σφόδρα εὐτελές. … καὶ σπανοπώγων τις ὀνομάζεται, δασυπώγων δὲ ἐν ταῖς Θεσμοφοριαζούσαις ; Ar.~\textit{Ach.}~119–121 ὦ θερμόβουλον πρωκτὸν ἐξυρημένε, / τοιόνδε γ᾽ ὦ πίθηκε τὸν πώγων᾽ ἔχων / εὐνοῦχος ἡμῖν ἦλθες ἐσκευασμένος ; Ar.~\textit{Nub.}~348–349 κομήτην / ἄγριόν τινα τῶν λασίων τούτων ; Ar.~\textit{Thesm.}~218–219 Ἀγάθων σὺ μέντοι ξυροφορεῖς ἑκάστοτε, / χρῆσόν τί νυν ἡμῖν ξυρόν ; Soph.~\textit{Trach.}~13 ἐκ δὲ δασκίου γενειάδος ; Archil. fr.~58~Bergk οὐ φιλέω μέγαν στρατηγὸν οὐδὲ διαπεπλιγμένον, / οὐδὲ βοστρύχοισι γαῦρον οὐδ᾽ ὑπεξυρημένον ; Theoc.~7.15 ἐκ μὲν γὰρ λασίοιο δασύτριχος εἶχε τράγοιο ; Plu.~Mor.~177a–b τοῦ δὲ Εὐριπίδου τὸν καλὸν Ἀγάθωνα περιλαμβάνοντος ἐν τῷ συμποσίῳ καὶ καταφιλοῦντος ἤδη γενειῶντα ; Ps.-Polem.~\textit{Phgn.}~70 λεῖος τὰ περὶ στῆθος καὶ βουβῶνα, δασυπώγων ; \textit{Vit.\,Eur.} (Σ Eur. vet., vita~2.1) ἐλέγετο δὲ καὶ βαθὺν πώγωνα θρέψαι καὶ ἐπὶ τῆς ὄψεως φακοὺς ἐσχηκέναι ; Σ \textit{Thesm.}~33 διαβάλλει τὸν Ἀγάθωνα πάλιν ὡς λεῖον καὶ ἀποτίλλοντα αὐτοῦ τὰς τρίχας τοῦ πώγωνος ; Hsch.~δ~305 δασύ· τραχύ. πυκνόν ; Hsch.~π~4488 πώγων· γένειον. ἄλλοι τὴν ἀναφορὰν τῆς φλογός ; Tosi~2183 Barba non facit philosophum — ἀπὸ πώγωνος φιλόσοφοι / σοφοί / σοφισταί}
\lp{34 μὰ τὸν Δί᾽ οὔτοι γ᾽]}{Ar.~\textit{Thesm.}~533 οὔ τοι μὰ τὴν Ἄγλαυρον ὦ γυναῖκες εὖ φρονεῖτε ; Ar.~\textit{Eccl.}~522 οὔτοι παρὰ τοῦ μοιχοῦ γε φήσεις ; Soph.~\textit{El.}~773 οὔτοι μάτην γε}
\lp{34 ὥστε κἀμέ γ᾽ εἰδέναι]}{Ar.~\textit{Nub.}~1252 οὐκ ἄρ᾽ ἀποδώσεις; — οὐχ ὅσον γέ μ᾽ εἰδέναι ; Ar.~\textit{Eccl.}~350 οὔκουν πονηρά γ᾽ ἐστὶν ὅ τι κἄμ᾽ εἰδέναι ; Plaut.~\textit{Epid.}~638 Non me nosti? — Quod quidem nunc veniat in mentem mihi ; Ter.~\textit{Hec.}~863 Nam nunquam ante hunc diem meis oculis eam, quod nossem, videram ; Pl.~\textit{Tht.}~145a ἦ οὖν ζωγραφικὸς Θεόδωρος; — οὔχ, ὅσον γέ με εἰδέναι ; Suid.~ο~734 Ὅ τι μ᾽ εἰδέναι: ἀντὶ τοῦ ὅσον ἐμὲ εἰδέναι. οὔκουν πονηρά γ᾽ ἐστιν, ὅ τι κἄμ᾽ εἰδέναι}
\lp{35 καὶ μὴν †βεβίνηκας† σύ γ᾽]}{Ar.~\textit{Thesm.}~200–201 καὶ μὴν σύ γ᾽ ὦ κατάπυγον εὐρύπρωκτος εἶ / οὐ τοῖς λόγοισιν ἀλλὰ τοῖς παθήμασιν}
\lp{35 †βεβίνηκας†]}{Ar.~\textit{Lys.}~1091–1092 ὥστ᾽ εἴ τις ἡμᾶς μὴ διαλλάξει ταχύ, / οὐκ ἔσθ᾽ ὅπως οὐ Κλεισθένη βινήσομεν ; Ar.~\textit{Thesm.}~50 μέλλει γὰρ ὁ καλλιεπὴς Ἀγάθων / πρόμος ἡμέτερος — μῶν βινεῖσθαι ; Ar.~\textit{Thesm.}~97–98 ἀλλ᾽ ἢ τυφλὸς μέν εἰμ᾽· ἐγὼ γὰρ οὐχ ὁρῶ / ἄνδρ᾽ οὐδέν᾽ ἐνθάδ᾽ ὄντα, Κυρήνην δ᾽ ὁρῶ ; Ar.~\textit{Thesm.}~204–206 δοκῶν γυναικῶν ἔργα νυκτερείσια … νὴ Δία βινεῖσθαι μὲν οὖν ; Pl.~Epigr. (Diog. Laert.~3.32) τὴν ψυχὴν Ἀγάθωνα φιλῶν ἐπὶ χείλεσιν εἶχον· / ἦλθε γὰρ ἡ τλήμων ὡς διαβησομένη ; Pl.~\textit{Smp.}~193b–c καὶ μή μοι ὑπολάβῃ Ἐρυξίμαχος, κωμῳδῶν τὸν λόγον, ὡς Παυσανίαν καὶ Ἀγάθωνα λέγω ; X.~\textit{Smp.}~8.32 Παυσανίας γε ὁ Ἀγάθωνος τοῦ ποιητοῦ ἐραστής ; Ael.~\textit{VH}~2.21 Ἀγάθωνος ἤρα τοῦ ποιητοῦ Παυσανίας ὁ ἐκ Κεραμέων. καὶ τοῦτο μὲν διατεθρύληται ; Ael.~\textit{VH}~13.4 Ἀγάθωνα τὸν τῆς τραγῳδίας ποιητὴν περιλαβὼν κατεφίλει ; Hsch.~β~466 βινεῖν· παρὰ Σόλωνι τὸ βίᾳ μίγνυσθαι. τὸ δὲ κατὰ νόμον ὀπύειν ; Hsch.~β~467 βινητιᾶν· τὸ πασχητιᾶν. καὶ τὸ ὄρεξιν ἐπὶ συνουσίαν ἔχειν ; Suid.~β~287–288 Βινεῖν: περαίνειν … οὐκ ἔσθ᾽ ὅπως μὴ Κλεισθένη βινήσομεν. οὗτος ἐπὶ θηλύτητι ἐκωμῳδεῖτο}
\lp{35 ἀλλ᾽ οὐκ οἶσθ᾽ ἴσως]}{Pherecr. fr.~7~Kock μὴ θαυμάσῃς· τῶν γὰρ προτενθῶν ἐσμέν, ἀλλ᾽ οὐκ οἶσθα σύ}
\lp{36 ἀλλ᾽ ἐκποδὼν πτήξωμεν]}{\textit{Test.} Suid.~π~3024 Πτήξωμεν, κρυβῶμεν. κυρίως δὲ πτήσσειν ἐστὶ τὸ φεύγειν καὶ φοβεῖσθαι. ἀλλ᾽ ἐκποδὼν πτήξωμεν. Ἀριστοφάνης}
\lp{36 ἀλλ᾽ ἐκποδὼν … ὡς ἐξέρχεται]}{Ar.~\textit{Ach.}~239–241 ἀλλὰ δεῦρο πᾶς / ἐκποδών· θύσων γὰρ ἁνὴρ ὡς ἔοικ᾽ ἐξέρχεται. — εὐφημεῖτε, εὐφημεῖτε ; Ar.~\textit{Pax}~1264–1267 χωρῶμεν ὦ τᾶν ἐκποδών. — νὴ τὸν Δί᾽, ὡς τὰ παιδί᾽ ἤδη ᾽ξέρχεται ; Eur.~\textit{El.}~107–111 εἰσορῶ γὰρ τήνδε προσπόλον τινά … ἑζώμεσθα κἀκπυθώμεθα ; Eur.~\textit{Hec.}~52–53 γεραιᾷ δ᾽ ἐκποδὼν χωρήσομαι / Ἑκάβῃ· περᾷ γὰρ ἥδ᾽ ὑπὸ σκηνῆς πόδα ; Eur.~\textit{Hel.}~1023 ἐγὼ δ᾽ ἀποστᾶσ᾽ ἐκποδὼν σιγήσομαι ; Aesch.~\textit{Cho.}~20–21 Πυλάδη, σταθῶμεν ἐκποδών, ὡς ἂν σαφῶς / μάθω γυναικῶν ἥτις ἥδε προστροπή ; Soph.~\textit{OC}~113–114 σιγήσομαί τε καὶ σύ μ᾽ ἐξ ὁδοῦ πόδα / κρύψον κατ᾽ ἄλσος}
\lp{36–41 ἐκποδών … ὡς ἐξέρχεται, θεράπων … μυρρίνας, εὔφημος … στόμα]}{Eur.~\textit{Ion}~78–79,~100 ὁρῶ γὰρ ἐκβαίνοντα Λοξίου γόνον / τόνδ᾽, ὡς πρὸ ναοῦ λαμπρὰ θῇ πυλώματα … στόμα τ᾽ εὔφημον φρουρεῖν ἀγαθόν}
\lp{36 πτήξωμεν]}{Ar.~\textit{Vesp.}~1490 πτήσσει Φρύνιχος ὥς τις ἀλέκτωρ — ; Ar.~\textit{Av.}~777 πτῆξε δὲ φῦλά τε ποικίλα θηρῶν ; Ar.~\textit{Lys.}~770 ἀλλ᾽ ὁπόταν πτήξωσι χελιδόνες εἰς ἕνα χῶρον ; Ar.~\textit{Ran.}~315 ἀλλ᾽ ἠρεμὶ πτήξαντες ἀκροασώμεθα ; Eur.~\textit{Andr.}~753 πτήξαντες οἵδε πρὸς βίαν ἄγωσί με ; Eur.~\textit{Ba.}~1035 οὐκέτι γὰρ δεσμῶν ὑπὸ φόβῳ πτήσσω ; Phryn. Trag. fr.~17~N.² ἔπτηξ᾽ ἀλέκτωρ δοῦλον ὣς κλίνας πτερόν ; \textit{Od.}~14.474 ἂν δόνακας καὶ ἕλος, ὑπὸ τεύχεσι πεπτηῶτες ; Σ \textit{Thesm.}~36 κρυβῶμεν. κυρίως δὲ πτήσσειν ἐστὶ τὸ φεύγειν καὶ φοβεῖσθαι ; Hsch.~ε~5565 ἔπτηξαν· ἐκρύβησαν. ἔπεσαν. ἐφοβήθησαν ; Hsch.~π~4220 πτῆξαι· δειλιᾶσαι ; Hsch.~π~4224 πτήσσων· ὑποπίπτων. κρυβόμενος}
\lp{36 ἐκποδών]}{Hsch.~ε~1635 ἐκποδών· ἐκ μέσου ; Suid.~ε~586 Ἐκποδών· γενικῇ. τὸ πρὸς ὀλίγον χρόνον ὑποστείλασθαι ; Zen.~1.90 Ἀνδρὸς κακῶς πράσσοντος ἐκποδὼν φίλοι}
\lp{36 ὡς ἐξέρχεται]}{Ar.~\textit{Thesm.}~70 τί οὖν ἐγὼ δρῶ; — περίμεν᾽, ὡς ἐξέρχεται ; Ar.~\textit{Thesm.}~95 σίγα. — τί δ᾽ ἔστιν; — Ἁγάθων ἐξέρχεται ; Ar.~\textit{Thesm.}~293 σὺ δ᾽ ἄπιθ᾽ ὦ Θρᾷττ᾽ ἐκποδών}
\lp{37 θεράπων τις αὐτοῦ]}{Ar.~\textit{Av.}~908–909 Μουσάων θεράπων ὀτρηρός ; Hes.~\textit{Th.}~99–100 Μουσάων θεράπων κλέεα προτέρων ἀνθρώπων}
\lp{37 πῦρ ἔχων καὶ μυρρίνας]}{Ar.~\textit{Vesp.}~860–862 πῦρ τις ἐξενεγκάτω / καὶ μυρρίνας καὶ τὸν λιβανωτὸν ἔνδοθεν, / ὅπως ἂν εὐξώμεσθα πρῶτα τοῖς θεοῖς ; Ar.~\textit{Av.}~43 κανοῦν δ᾽ ἔχοντε καὶ χύτραν καὶ μυρρίνας ; Ar.~\textit{Ran.}~871–872 ἴθι νυν λιβανωτὸν δεῦρό τις καὶ πῦρ δότω, / ὅπως ἂν εὔξωμαι πρὸ τῶν σοφισμάτων ; Eur.~\textit{Alc.}~172 πτόρθων ἀποσχίζουσα μυρσίνης φόβην ; Eur.~\textit{El.}~777–778 κυρεῖ δὲ κήποις ἐν καταρρύτοις βεβώς, / δρέπων τερείνης μυρσίνης κάρᾳ πλόκους ; Eur.~\textit{Ion}~120 μυρσίνας ἱερὰν φόβαν ; Hdt.~7.54.1 θυμιήματά τε παντοῖα ἐπὶ τῶν γεφυρέων καταγίζοντες καὶ μυρσίνῃσι στορνύντες τὴν ὁδόν ; Poll.~1.27–28 στεφανώματα προσφέρειν, μυῤῥίνην στεφανωτρίδα … θαλλούς, μυῤῥίνας, ἄνθη ; Zen.~5.34 Οὐδὲ πυρφόρος ἐλείφθη: ἐπὶ τῶν μεγάλων κινδύνων εἴρηται ἡ παροιμία}
\lp{37 μυρρίνας]}{Ar.~\textit{Eq.}~964 ψωλὸν γενέσθαι δεῖ σε μέχρι τοῦ μυρρίνου ; Ar.~\textit{Nub.}~1364–1365 ἔπειτα δ᾽ ἐκέλευσ᾽ αὐτὸν ἀλλὰ μυρρίνην λαβόντα / τῶν Αἰσχύλου λέξαι τί μοι ; Ar.~\textit{Thesm.}~448 στεφανηπλοκοῦσ᾽ ἔβοσκον ἐν ταῖς μυρρίναις ; Σ Ar.~\textit{Eq.}~964 μύρτον δὲ πεποικιλμένον ἐκάλουν οἱ παλαιοὶ τὸ γυναικεῖον αἰδοῖον ; Hsch.~μ~1915 μυρρίνης κλάδος, ἢ δάφνης· παρὰ πότον μυρρίνης ⟨κλάδον⟩ ἦν σύνηθες διδόναι τοῖς κατακειμένοις ἐκ διαδοχῆς ὑπὲρ τοῦ ᾆσαι, ἀντὶ τοῦ βαρβίτου ; Hsch.~μ~1917 μύρρινος· ἐπειδὴ μύρρινον στέφανον περιετίθεντο ; Hsch.~μ~1926 μύρτος· ἡ μυρρίνη. καὶ τὸ γυναικεῖον αἰδοῖον· οἱ δὲ τὸ σχίσμα τῆς γυναικός ; Poll.~6.108 καὶ μυρρίνην ἐπὶ δεξιὰ περιφέροντές τινες καὶ ἔκπωμα καὶ λύραν ᾄδειν ἠξίουν ; Zen.~1.19 Ἄιδειν πρὸς μυῤῥίνην: ἔθος ἦν τὸν μὴ δυνάμενον ἐν τοῖς συμποσίοις ᾆσαι, δάφνης κλῶνα ἢ μυῤῥίνης λαβόντα πρὸς τοῦτον ᾄδειν ; Hsch.~α~1762 ἄιδειν πρὸς μυρρίνην· ἔθος ἐν τοῖς συμποσίοις τὸν μὴ δυνάμενον ᾆσαι δάφνης κλῶνα ἢ μυρρίνης λαβόντα πρὸς αὐτὴν ἄιδειν}
\lp{37–41 μυρρίνας … θίασος]}{Ar.~\textit{Ran.}~155–156 ὄψει τε φῶς κάλλιστον ὥσπερ ἐνθάδε, / καὶ μυρρινῶνας καὶ θιάσους εὐδαίμονας}
\lp{38 προθυσόμενος]}{Ar.~\textit{Ran.}~871–874 ὅπως ἂν εὔξωμαι πρὸ τῶν σοφισμάτων ; Ar.~\textit{Plut.}~660–661 ἐπεὶ δὲ βωμῷ πόπανα καὶ προθύματα / καθωσιώθη πέλανος Ἡφαίστου φλογί ; Eur.~\textit{IA}~1310–1311 πρόθυμα δ᾽ / ἔλαβεν Ἄρτεμις πρὸς Ἴλιον ; Eur.~\textit{Ion}~805 παιδὸς προθύσων ξένια καὶ γενέθλια ; Eur.~\textit{Supp.}~28–29 τυγχάνω δ᾽ ὑπὲρ χθονὸς / ἀρότου προθύουσ᾽ ; Pl.~\textit{Cra.}~401d τὸ γὰρ πρὸ πάντων θεῶν τῇ Ἑστίᾳ πρώτῃ προθύειν ; Plu.~\textit{Lyc.}~21.4 καὶ γὰρ ἐν ταῖς μάχαις προεθύετο ταῖς Μούσαις ὁ βασιλεύς ; Σ \textit{Thesm.}~38 ὅτε γὰρ δρᾶμα ποιῆσαι ἤθελον, πρότερον θυσίας ἐποίουν. — λείπει τὸ ὡς, ὡς ἔοικε ; Σ Ar.~\textit{Plut.}~660 σημαίνει δὲ τὰ προκατάργματα, ἢ τὰ πρὸ τῆς θυσίας γινόμενα θυμιάματα ; Suid.~π~2448 Προθύματα: αἱ ὄλυραι. παρὰ τὸ προθύεσθαι τῶν ἱερείων, ἢ αἱ κριθαί, ἢ ὁ λιβανωτός ; Zen.~1.40 Ἀφ᾽ Ἑστίας ἀρχόμενοι: μετενήνεκται δὲ ἀπὸ τῶν περὶ τὰ ἱερὰ δρωμένων}
\lp{38 ἔοικε]}{Ar.~\textit{Plut.}~1097–1098 οὐδεὶς ἔοικεν· ἀλλὰ δῆτα τὸ θύριον / φθεγγόμενον ἄλλως κλαυσιᾷ}
\lp{38 τῆς ποιήσεως]}{Ar.~\textit{Ran.}~868 ὅτι ἡ ποίησις οὐχὶ συντέθνηκέ μοι ; Ar.~\textit{Ran.}~1366 ὅπερ ἐξελέγξει τὴν ποίησιν νῷν μόνον ; Hdt.~2.82.1 τῶν Ἑλλήνων οἱ ἐν ποιήσι γενόμενοι ; Pl.~\textit{Smp.}~205b–c οἶσθ᾽ ὅτι ποίησίς ἐστί τι πολύ}
\lp{39–41 εὔφημος … μελάθρων]}{\textit{Test.} Suid.~ε~3792 Ἀριστοφάνης· εὔφημος πᾶς ἔστω λεὼς στόμα συγκλείσας. ἐνδημεῖ γὰρ θίασος Μουσῶν ἔνδον μελάθρων ; Eur.~\textit{Ba.}~68–70 τίς ὁδῷ τίς ὁδῷ; τίς; / μελάθροις ἔκτοπος ἔστω, στόμα τ᾽ εὔφη- / μον ἅπας ἐξοσιούσθω ; Hor.~\textit{Carm.}~3.1.2–4 favete linguis: carmina non prius / audita Musarum sacerdos / virginibus puerisque canto}
\lp{39–42 εὔφημος … μελοποιῶν]}{Ar.~\textit{Ach.}~395–399 οὐκ ἔνδον ἔνδον ἐστίν, εἰ γνώμην ἔχεις … αὐτὸς δ᾽ ἔνδον ἀναβάδην ποιεῖ / τραγῳδίαν ; Eur.~fr.~773,~58–73~N.² (Phaéthon) κόσμον δ᾽ ὑμεναίων δεσποσύνων … ἔχειν χρὴ / στόμ᾽ ἐν ἡσυχίᾳ· … Ὠκεανοῦ πεδίων οἰκήτορες, εὐφαμεῖτ᾽, ὦ, / ἐκτόπιοί τε δόμων ἀπαείρετε, / ὦ ἴτε, λαοί. … ἀλλὰ σῖγ᾽ ἔστω λεώς}
\lp{39 εὔφημος πᾶς ἔστω λαός]}{Ar.~\textit{Ach.}~237–238 εὐφημεῖτε, εὐφημεῖτε. — σῖγα πᾶς. ἠκούσατ᾽ ἄνδρες ἆρα τῆς εὐφημίας ; Ar.~\textit{Nub.}~263 εὐφημεῖν χρὴ τὸν πρεσβύτην καὶ τῆς εὐχῆς ἐπακούειν ; Ar.~\textit{Pax}~96–101 εὐφημεῖν χρὴ καὶ μὴ φλαῦρον / μηδὲν γρύζειν … καὶ τοὺς πρωκτοὺς ἐπικλῄειν ; Ar.~\textit{Av.}~1718–1719 ἀλλὰ χρὴ θεᾶς / Μούσης ἀνοίγειν ἱερὸν εὔφημον στόμα ; Ar.~\textit{Thesm.}~295 εὐφημία ᾽στω, εὐφημία ᾽στω ; Ar.~\textit{Ran.}~354–356 εὐφημεῖν χρὴ κἀξίστασθαι τοῖς ἡμετέροισι χοροῖσιν … ἢ γενναίων ὄργια Μουσῶν μήτ᾽ εἶδεν μήτ᾽ ἐχόρευσεν ; Eur.~\textit{Hec.}~532 Σιγᾶτ᾽, Ἀχαιοί, σῖγα πᾶς ἔστω λεώς ; Eur.~\textit{Hipp.}~724 εὔφημος ἴσθι ; Eur.~\textit{Or.}~1327 εὔφημος ἴσθι· τί δὲ νεώτερον λέγεις ; Eur.~\textit{Rhes.}~730 σῖγα πᾶς ὕφιζ᾽· ἴσως γὰρ ἐς βόλον τις ἔρχεται ; Call.~\textit{Ap.}~17 εὐφημεῖτ᾽ ἀίοντες ἐπ᾽ Ἀπόλλωνος ἀοιδῇ ; Mesom.~\textit{Hymn.\,Sol.}~1–6 εὐφαμείτω πᾶς αἰθήρ … μέλλει δὲ πρὸς ἡμᾶς βαίνειν / Φοῖβος ; Σ \textit{Thesm.}~39 παρὰ τοῖς τραγικοῖς οὕτω λέγεται, πᾶς ἔστω λαός ; Σ Eur.~\textit{Hec.}~532 σῖγα πᾶς ἔστω λαός: μιμεῖται τὸ πολλάκις κηρύσσειν τὰ αὐτὰ τοὺς κήρυκας ; Hsch.~ε~7264 εὐφημίαν· σιωπήν. εὐλογίαν ; Hsch.~λ~792 λεώς· Ἀττικῶς ὄχλος, λαός ; Tosi~2083 Favete linguis!}
\lp{39–40 εὔφημος … στόμα συγκλῄσας]}{Ar.~\textit{Eq.}~1316 εὐφημεῖν χρὴ καὶ στόμα κλῄειν καὶ μαρτυριῶν ἀπέχεσθαι ; Aesch.~\textit{Ag.}~1247 εὔφημον, ὦ τάλαινα, κοίμησον στόμα}
\lp{40–42 ἐπιδημεῖ γὰρ θίασος Μουσῶν … μελοποιῶν]}{\textit{Test.} Suid.~θ~378 Ἀριστοφάνης· ἐπιδημεῖ γὰρ θίασος Μουσῶν ἔνδον μελάθρων τῶν δεσποσύνων μελοποιῶν}
\lp{40 στόμα συγκλῄσας]}{Ar.~\textit{Ran.}~837–838 ἄνθρωπον ἀγριοποιὸν αὐθαδόστομον, / ἔχοντ᾽ ἀχάλινον ἀκρατὲς ἀπύλωτον στόμα ; Eur.~\textit{Hec.}~664 οὐ ῥᾴδιον βροτοῖσιν εὐφημεῖν στόμα ; Eur.~\textit{Hipp.}~498 ὦ δεινὰ λέξασ᾽, οὐχὶ συγκλῄσεις στόμα ; Soph.~\textit{OC}~1051–1053 ὧν καὶ χρυσέα / κλῂς ἐπὶ γλώσσᾳ βέβακε / προσπόλων Εὐμολπιδᾶν ; Phot.~p.~145~Porson Κλῆϊσαι: οἱ ἀρχαῖοι λέγουσιν· οὐ κλεῖσαι· καὶ κλῆϊδα· οὕτω καὶ οἱ τραγικοὶ καὶ Θουκυδίδης ; Zen.~2.70 Βοῦς ἐπὶ γλώττης: παροιμία ἐπὶ τῶν μὴ δυναμένων παρρησιάζεσθαι}
\lp{40 ἐπιδημεῖ]}{Ar.~fr.~348~K.-A. μήτε Μούσας ἀνακαλεῖν ἑλικοβοστρύχους … ἐνθάδε γάρ εἰσιν, ὥς φησιν ὁ διδάσκαλος ; Call.~\textit{Ap.}~12–13 μήτε σιωπηλὴν κίθαριν μήτ᾽ ἄψοφον ἴχνος / τοῦ Φοίβου τοὺς παῖδας ἔχειν ἐπιδημήσαντος ; Pl.~\textit{Prt.}~309d Πρωταγόρας ἐπιδεδήμηκεν; — τρίτην γε ἤδη ἡμέραν ; X.~\textit{Mem.}~1.2.61 Λίχας μὲν γὰρ ταῖς γυμνοπαιδίαις τοὺς ἐπιδημοῦντας ἐν Λακεδαίμονι ξένους ἐδείπνιζε ; Luc.~\textit{Symp.}~7 καὶ ὅλως θεοῦ ἐπιδημία τὸ πρᾶγμα ἦν Ἴων ὁ θαυμαστὸς συμπαρών ; Hsch.~ε~2777 ἐνδημία· παρουσία ; Hsch.~ε~4722 ἐπιδημία· παρουσία}
\lp{41 θίασος Μουσῶν]}{Ar.~\textit{Nub.}~297 ἀλλ᾽ εὐφήμει· μέγα γάρ τι θεῶν κινεῖται σμῆνος ἀοιδαῖς ; Ar.~\textit{Ran.}~92–93 ἐπιφυλλίδες ταῦτ᾽ ἐστὶ καὶ στωμύλματα, / χελιδόνων μουσεῖα, λωβηταὶ τέχνης ; Eur.~\textit{Ba.}~55–56 ἀλλ᾽, ὦ λιποῦσαι Τμῶλον ἔρυμα Λυδίας, / θίασος ἐμός, γυναῖκες ; Eur.~\textit{IA}~1059–1060 θίασος ἔμολεν ἱπποβάτας / Κενταύρων ἐπὶ δαῖτα ; Eur.~\textit{IT}~1146 ματρὸς ἡλίκων θιάσους ; Eur.~\textit{Phoen.}~796 ἀσπιδοφέρμονα θίασον ἐνόπλιον ; D.~18.260 ἐν δὲ ταῖς ἡμέραις τοὺς καλοὺς θιάσους ἄγων διὰ τῶν ὁδῶν ; [Plu.] \textit{De\,mus.}~2 ἄγε δή, ὦ μουσικῆς θιασῶται ; Σ \textit{Thesm.}~41 ὁ ἱερὸς χορός ; Hsch.~θ~572 θίασον· εὐωχίαν. καὶ πλῆθος, οὐ μόνον τὸ Βακχικόν, ἀλλὰ ⟨καὶ⟩ τὸ ἐρανικόν ; Hsch.~θ~573 θίασος· χοροῦ σύστασις, ἢ σύλλογος. καὶ τὸ ψιλὸν πλῆθος· ἔσθ᾽ ὅτε δὲ καὶ τὸ Βακχικὸν ἢ ἑσμὸς γυναικῶν ; Hsch.~μ~1742 μοῦσα· τέχνη. καὶ σοφοὺς καὶ μουσικοὺς τοὺς τεχνίτας ἔλεγον ; Suid.~θ~379–380 Θίασος: ἱερὸς χορός … Ἴων δὲ ἐπὶ παντὸς ἀθροίσματος τάττει τοὔνομα ; Suid.~θ~381 Θιασώτης Κότυος. Κότυς, δαίμων παρὰ Κορινθίοις τιμώμενος, ἔφορος τῶν αἰσχρῶν. ὅθεν οἱ Κλεισθένεις καὶ Τίμαρχοι ; Zen.~1.89 Ἀνεῳγμέναι Μουσῶν θύραι: ἐπὶ τῶν ἐξ ἑτοίμου λαμβανόντων τὰ κάλλιστα τῶν ἐν παιδείᾳ ; Zen.~4.27 Ὁ Θάμυρις πολλῶν ἐπὶ κάλλει διενεγκὼν, ἤρξατο πρῶτον ἐρᾶν ἀῤῥένων. Ἀσκήσας δὲ κιθαρῳδίαν, ταῖς Μούσαις ἤρισε ; Diogenian.~6.50 Μετὰ Μουσῶν κόρυδος: ἐπὶ τῶν ἀπαιδεύτων πεπαιδευμένοις συνόντων. Ὁμοία, Καρικὴ Μοῦσα ; Σ Ar.~\textit{Ran.}~93 χελιδόνων μουσεῖα· παρὰ τὰ ἐν Ἀλκμήνῃ Εὐριπίδου· πολὺς δ᾽ ἀνεῖρπε κισσός, εὐφυὴς κλάδος, χελιδόνων μουσεῖον. … τὸ χελιδόνων μουσεῖα παροιμία ἐπὶ τῶν πολυλόγων καὶ ἐπαχθῶν ταττομένη}
\lp{41–42 θίασος Μουσῶν … μελοποιῶν]}{Pl.~\textit{Ion}~533e–534b οἱ μελοποιοὶ οὐκ ἔμφρονες ὄντες τὰ καλὰ μέλη ταῦτα ποιοῦσιν … βακχεύουσι καὶ κατεχόμενοι}
\lp{41 ἔνδον μελάθρων]}{Ar.~\textit{Thesm.}~874 Πρωτέως τάδ᾽ ἐστὶ μέλαθρα ; Eur.~\textit{Hec.}~1100–1101 ἀμπτάμενος οὐράνιον / ὑψιπετὲς ἐς μέλαθρον ; Hsch.~μ~623 μέλαθρα· οἰκίαι. ὑπέρθυρα. δοκὸς δὲ καὶ ἡ διάτονος}
\lp{42 τῶν δεσποσύνων]}{Anaxandr. fr.~41.33~Kock (=~42.33~K.-A.) τούτων δ᾽ ἔσται πολὺ σεμνότερον καὶ λαμπρότερον / παρὰ δεσποσύνοις τοῖς ἡμετέροις ; Eur.~\textit{Hec.}~99 τὰς δεσποσύνους σκηνὰς προλιποῦσ᾽ ; Aesch.~\textit{Cho.}~942 ἐπολολύξατ᾽ ὦ δεσποσύνων δόμων ; Tyrt. fr.~6.2~Bergk δεσποσύνοισι φέροντες ἀναγκαίης ὑπὸ λυγρῆς ; Tim. Pers.~136 (= PMG~791.125) Ματρὸς οὐρείας δεσπόσυνα γόνατα πεσεῖν}
\lp{42 μελοποιῶν]}{Ar.~\textit{Thesm.}~67 καὶ γὰρ μελοποιεῖν ἄρχεται· χειμῶνος οὖν / ὄντος κατακάμπτειν τὰς στροφὰς οὐ ῥᾴδιον ; Ar.~\textit{Ran.}~1250 μελοποιὸν ὄντα καὶ ποιοῦντα ταὔτ᾽ ἀεί ; Ar.~\textit{Ran.}~1327–1328 ἀνὰ τὸ δωδεκαμήχανον / Κυρήνης μελοποιῶν ; Pl.~\textit{Prt.}~326a ἐπειδὰν κιθαρίζειν μάθωσιν, ἄλλων αὖ ποιητῶν ἀγαθῶν ποιήματα διδάσκουσι μελοποιῶν ; Hsch.~π~2891 ἦν δὲ Κολοφώνιος μελοποιὸς ὁ Πολύμνηστος, εὐμελὴς πάνυ}
